% !TeX root = ../main.tex
% chapters/appendix.tex

\appendix
\label{chap:appendix} % 付録セクション全体の開始点

% =============================================================
% 1. プログラムと構成（ 先頭に配置 ）
% =============================================================
\input{chapters/sections/appendix/programs_overview.tex}

% =============================================================
% 2. 第 3 章（ モデル構築 ）で参照される基礎的な導出
% =============================================================
\input{chapters/sections/appendix/time_consistency.tex}
% !TeX root = ../../main.tex
% sections/app/appendix_cost_minimization.tex

\chapter{費用最小化問題からの需要関数と価格指数の導出（第1段階）}
\label{chap:appendix_cost_minimization}

本付録の目的は、総名目生産額が、集計産出量と真の価格指数の単純な積で表せること、そしてその関係を担保するために集計産出量が必然的にCES集計の形でなければならないことを証明することである。

\section*{定理2：個別財への需要関数と真の価格指数}
\label{sec:appendix_cost_minimization_theorem2}

家計\(h\)が、ある時点\(t\)において、所与の量の国内財バスケット\(c_t^{h\to H}=\left[\sum_{h'\in H}(c_t^{h\to h'})^{\frac{\theta^H-1}{\theta^H}}\right]^{\frac{\theta^H}{\theta^H-1}}\)を、個別財の価格\(\{p_t^{h'}\}\)の下で費用最小化的に購入するとき、以下の関係が成立する。
\begin{enumerate}
    \item\label{itm:appendix_cost_minimization_theorem2_1}家計\(h\)の個別財\(h'\)への需要関数は、次式で与えられる。
    \begin{equation}
    c_t^{h\to h'}=\left(\frac{p_t^{h'}}{p_t^H}\right)^{-\theta^H}c_t^{h\to H}
    \label{eq:appendix_cost_minimization_theorem2_individual_demand}
    \end{equation}
    \item\label{itm:appendix_cost_minimization_theorem2_2}上記の需要関数に現れる\(p_t^H\)は、国内財バスケットの「真の価格指数」であり、個々の財価格から以下のように定義される。この値は全ての国内家計で共通である。
    \begin{equation}
    p_t^H=\left[\sum_{h'\in H}(p_t^{h'})^{1-\theta^H}\right]^{\frac{1}{1-\theta^H}}
    \label{eq:appendix_cost_minimization_theorem2_price_index}
    \end{equation}
\end{enumerate}

\section*{証明}
\label{sec:appendix_cost_minimization_proof}

この定理を、自国家計\(h\)の場合について証明する。外国の家計\(f\)についても全く対称的な手順で証明可能である。

\subsection*{1.費用最小化問題の設定}
\label{sec:appendix_cost_minimization_sub1}
家計が直面する問題は、時点\(t\)において、所与の国内財バスケット\(c_t^{h\to H}\)を達成するための総費用を最小化することである。
\begin{itemize}
    \item\label{itm:appendix_cost_minimization_sub1_obj} \textbf{最小化対象}:
    \begin{equation}
    \min_{\{c_t^{h\to h'}\}_{h'\in H}}\sum_{h'\in H}p_t^{h'}c_t^{h\to h'}
    \label{eq:appendix_cost_minimization_sub1_objective}
    \end{equation}
    \item\label{itm:appendix_cost_minimization_sub1_const} \textbf{制約条件}:
    \begin{equation}
    c_t^{h\to H}=\left[\sum_{h'\in H}(c_t^{h\to h'})^{\frac{\theta^H-1}{\theta^H}}\right]^{\frac{\theta^H}{\theta^H-1}}
    \label{eq:appendix_cost_minimization_sub1_constraint}
    \end{equation}
\end{itemize}

\subsection*{2.ラグランジュ関数と一階の条件（FOC）}
\label{sec:appendix_cost_minimization_sub2}
この問題を解くために、ラグランジュ関数\(\mathcal{L}^{h\to H}\)を設定する。ラグランジュ乗数を\(\mu_t^{h\to H}\)とする。
\begin{equation}
\mathcal{L}^{h\to H}=\sum_{h'\in H}p_t^{h'}c_t^{h\to h'}-\mu_t^{h\to H}\left(\left[\sum_{h'\in H}(c_t^{h\to h'})^{\frac{\theta^H-1}{\theta^H}}\right]^{\frac{\theta^H}{\theta^H-1}}-c_t^{h\to H}\right)
\label{eq:appendix_cost_minimization_sub2_lagrangian}
\end{equation}
このラグランジュ関数を、任意の個別財\(c_t^{h\to h'}\)で偏微分し、ゼロと置くことで一階の条件（FOC）が得られる。
\begin{align}
\frac{\partial\mathcal{L}^{h\to H}}{\partial c_t^{h\to h'}}=p_t^{h'}-\mu_t^{h\to H}\cdot\frac{\theta^H}{\theta^H-1}\left[\sum_{i\in H}(c_t^{h\to i})^{\frac{\theta^H-1}{\theta^H}}\right]^{\frac{\theta^H}{\theta^H-1}-1}\cdot\frac{\theta^H-1}{\theta^H}(c_t^{h\to h'})^{\frac{\theta^H-1}{\theta^H}-1}&=0 \label{eq:appendix_cost_minimization_sub2_foc_step1} \\
p_t^{h'}&=\mu_t^{h\to H}\cdot(c_t^{h\to H})^{\frac{1}{\theta^H}}\cdot(c_t^{h\to h'})^{-\frac{1}{\theta^H}} \label{eq:appendix_cost_minimization_sub2_foc_step2} \\
p_t^{h'}&=\mu_t^{h\to H}\left(\frac{c_t^{h\to H}}{c_t^{h\to h'}}\right)^{\frac{1}{\theta^H}}
\label{eq:appendix_cost_minimization_sub2_foc_final}
\end{align}

\subsection*{3.FOCからの需要関数の導出}
\label{sec:appendix_cost_minimization_sub3}
上記で得られたFOCを、個別財への需要量\(c_t^{h\to h'}\)について解く。
\begin{align}
    \frac{p_t^{h'}}{\mu_t^{h\to H}}&=\left(\frac{c_t^{h\to H}}{c_t^{h\to h'}}\right)^{\frac{1}{\theta^H}} \label{eq:appendix_cost_minimization_sub3_demand_step1} \\
    \left(\frac{p_t^{h'}}{\mu_t^{h\to H}}\right)^{\theta^H}&=\frac{c_t^{h\to H}}{c_t^{h\to h'}}
    \label{eq:appendix_cost_minimization_sub3_demand_step2}
\end{align}
これによりラグランジュ乗数\(\mu_t^{h\to H}\)を含む需要関数が導出される。
\begin{equation}
c_t^{h\to h'}=\left(\frac{p_t^{h'}}{\mu_t^{h\to H}}\right)^{-\theta^H}c_t^{h\to H}
\tag{\ref{eq:model_households_indices_aggregation_stage1a_home_c_h_H_problem_c_h_h_eq}}
\end{equation}

\subsection*{4.真の価格指数\(p_t^H\)の導出}
\label{sec:appendix_cost_minimization_sub4}
次に、ラグランジュ乗数\(\mu_t^{h\to H}\)の具体的な形を導出する。ステップ3で得た需要関数を、制約条件であるCES集計式に代入する。
\begin{equation}
c_t^{h\to H}=\left[\sum_{h'\in H}\left(\left(\frac{p_t^{h'}}{\mu_t^{h\to H}}\right)^{-\theta^H}c_t^{h\to H}\right)^{\frac{\theta^H-1}{\theta^H}}\right]^{\frac{\theta^H}{\theta^H-1}}
\label{eq:appendix_cost_minimization_sub4_ces_substitution}
\end{equation}
式を整理していく。
\begin{align}
    c_t^{h\to H}&=\left[\sum_{h'\in H}\left(\frac{p_t^{h'}}{\mu_t^{h\to H}}\right)^{1-\theta^H}(c_t^{h\to H})^{\frac{\theta^H-1}{\theta^H}}\right]^{\frac{\theta^H}{\theta^H-1}} \label{eq:appendix_cost_minimization_sub4_ces_step1} \\
    &=\left[(c_t^{h\to H})^{\frac{\theta^H-1}{\theta^H}}(\mu_t^{h\to H})^{\theta^H-1}\sum_{h'\in H}(p_t^{h'})^{1-\theta^H}\right]^{\frac{\theta^H}{\theta^H-1}} \label{eq:appendix_cost_minimization_sub4_ces_step2} \\
    &=(c_t^{h\to H})\cdot(\mu_t^{h\to H})^{\theta^H}\cdot\left[\sum_{h'\in H}(p_t^{h'})^{1-\theta^H}\right]^{\frac{\theta^H}{\theta^H-1}}
    \label{eq:appendix_cost_minimization_sub4_ces_final}
\end{align}
両辺の\(c_t^{h\to H}\)を消去し、\(\mu_t^{h\to H}\)について解く。
\begin{align}
    1&=(\mu_t^{h\to H})^{\theta^H}\left[\sum_{h'\in H}(p_t^{h'})^{1-\theta^H}\right]^{\frac{\theta^H}{\theta^H-1}} \label{eq:appendix_cost_minimization_sub4_mult_step1} \\
    (\mu_t^{h\to H})^{-\theta^H}&=\left[\sum_{h'\in H}(p_t^{h'})^{1-\theta^H}\right]^{\frac{\theta^H}{\theta^H-1}} \label{eq:appendix_cost_minimization_sub4_mult_step2} \\
    \mu_t^{h\to H}&=\left[\sum_{h'\in H}(p_t^{h'})^{1-\theta^H}\right]^{\frac{1}{1-\theta^H}}
    \label{eq:appendix_cost_minimization_sub4_mult_final}
\end{align}
このラグランジュ乗数\(\mu_t^{h\to H}\)は、個々の家計\(h\)に依存しない共通の価格リストのみで決定されるため、全ての国内家計で共通の値をとる。
本稿では、この家計が直面する真の価格指数を\(p_t^H\)と定義する。
\begin{equation}
p_t^H\equiv\mu_t^{h\to H}=\left[\sum_{h'\in H}(p_t^{h'})^{1-\theta^H}\right]^{\frac{1}{1-\theta^H}}
\label{eq:appendix_cost_minimization_sub4_price_index_def}
\end{equation}
この\(p_t^H\)をステップ3の需要関数に代入することで、定理の項目1が証明される（証明終）。      % 家計の需要関数
% !TeX root = ../../main.tex
% sections/app/appendix_cpi_derivation.tex

\chapter{総消費価格指数と財バスケットへの需要関数の導出（第2段階）}
\label{chap:appendix_cpi_derivation}

\section*{定理3：財バスケットへの需要関数と総消費価格指数}
\label{sec:appendix_cpi_derivation_theorem3}

家計\(h\)が、ある時点\(t\)において、所与の量の総消費バスケット\(c_t^{h\to W}\)を、国内財バスケットの真の価格指数\(p_t^H\)と外国財バスケットの真の価格指数\(p_t^{F*}\)の下で費用最小化的に購入するとき、以下の関係が成立する。
\begin{enumerate}
    \item\label{itm:appendix_cpi_derivation_theorem3_1}国内財バスケット\(c_t^{h\to H}\)および外国財バスケット\(c_t^{h\to F}\)への需要関数は、それぞれ次式で与えられる。
    \begin{equation}
    c_t^{h\to H}=\alpha^H\frac{p_t^{H\to W}}{p_t^H}c_t^{h\to W}
    \label{eq:appendix_cpi_derivation_theorem3_demand_domestic}
    \end{equation}
    \begin{equation}
    c_t^{h\to F}=(1-\alpha^H)\frac{p_t^{H\to W}}{e_tp_t^{F*}}c_t^{h\to W}
    \label{eq:appendix_cpi_derivation_theorem3_demand_foreign}
    \end{equation}
    \item\label{itm:appendix_cpi_derivation_theorem3_2}上記の需要関数に現れる\(p_t^{H\to W}\)は、総消費バスケットの価格指数（CPI）であり、各財バスケットの真の価格指数から以下のように定義される。
    \begin{equation}
    p_t^{H\to W}=(p_t^H)^{\alpha^H}(e_tp_t^{F*})^{1-\alpha^H}
    \label{eq:appendix_cpi_derivation_theorem3_cpi_def}
    \end{equation}
\end{enumerate}

\section*{証明}
\label{sec:appendix_cpi_derivation_proof}

この定理を、自国家計\(h\)の場合について証明する。

\subsection*{1.費用最小化問題の設定}
\label{sec:appendix_cpi_derivation_sub1}
この第2段階では、家計が時点\(t\)において、所与の総消費量\(c_t^{h\to W}\)を達成するために、国内財バスケット\(c_t^{h\to H}\)と海外財バスケット\(c_t^{h\to F}\)をどのように組み合わせれば総費用を最小化できるかを分析する。
\begin{itemize}
    \item\label{itm:appendix_cpi_derivation_sub1_obj}\textbf{最小化すべき総費用}:
    \begin{equation}
    \min_{c_t^{h\to H},c_t^{h\to F}}p_t^Hc_t^{h\to H}+e_tp_t^{F*}c_t^{h\to F}
    \label{eq:appendix_cpi_derivation_sub1_objective}
    \end{equation}
    \item\label{itm:appendix_cpi_derivation_sub1_const}\textbf{制約条件}（目標とする正規化されたバスケットの量）:
    \begin{equation}
    c_t^{h\to W}\equiv\frac{(c_t^{h\to H})^{\alpha^H}(c_t^{h\to F})^{1-\alpha^H}}{(\alpha^H)^{\alpha^H}(1-\alpha^H)^{1-\alpha^H}}
    \label{eq:appendix_cpi_derivation_sub1_constraint}
    \end{equation}
\end{itemize}

\subsection*{2.ラグランジュ関数と一階の条件（FOC）}
\label{sec:appendix_cpi_derivation_sub2}
制約式の対数を取ると扱いやすい。ラグランジュ関数\(\mathcal{L}^{h\to W}\)を設定する（この段階のラグランジュ乗数を\(\eta_t^h\)とする）。



\begin{equation}
\mathcal{L}^{h\to W}=p_t^Hc_t^{h\to H}+e_tp_t^{F*}c_t^{h\to F}-\eta_t^h\left(\alpha^H\ln c_t^{h\to H}+(1-\alpha^H)\ln c_t^{h\to F}-\ln c_t^{h\to W}-\text{const.}\right)
\label{eq:appendix_cpi_derivation_sub2_lagrangian}
\end{equation}
これを\(c_t^{h\to H}\)と\(c_t^{h\to F}\)でそれぞれ偏微分してゼロと置くと、以下の一階の条件（FOC）が得られる。
\begin{align}
\frac{\partial\mathcal{L}^{h\to W}}{\partial c_t^{h\to H}}=p_t^H-\eta_t^h\frac{\alpha^H}{c_t^{h\to H}}=0\quad&\implies\quad p_t^Hc_t^{h\to H}=\alpha^H\eta_t^h \label{eq:appendix_cpi_derivation_sub2_foc_h}\\
\frac{\partial\mathcal{L}^{h\to W}}{\partial c_t^{h\to F}}=e_tp_t^{F*}-\eta_t^h\frac{1-\alpha^H}{c_t^{h\to F}}=0\quad&\implies\quad e_tp_t^{F*}c_t^{h\to F}=(1-\alpha^H)\eta_t^h \label{eq:appendix_cpi_derivation_sub2_foc_f}
\end{align}

\subsection*{3.需要関数の導出（\(\eta_t^h\)を含む形）}
\label{sec:appendix_cpi_derivation_sub3}
上記の一階の条件をそれぞれ\(c_t^{h\to H}\)と\(c_t^{h\to F}\)について解くと、ラグランジュ乗数\(\eta_t^h\)を含む形で各バスケットへの需要関数が得られる。
\begin{align}
c_t^{h\to H}=\frac{\alpha^H\eta_t^h}{p_t^H}\quad,\quad c_t^{h\to F}=\frac{(1-\alpha^H)\eta_t^h}{e_tp_t^{F*}}
\label{eq:appendix_cpi_derivation_sub3_demand_with_multiplier}
\end{align}

\subsection*{4.ラグランジュ乗数\(\eta_t^h\)の導出}
\label{sec:appendix_cpi_derivation_sub4}
ステップ3で得た需要関数を、制約条件である総消費バスケットの定義式に代入する。
\begin{align}
c_t^{h\to W}&=\frac{1}{(\alpha^H)^{\alpha^H}(1-\alpha^H)^{1-\alpha^H}}\left(\frac{\alpha^H\eta_t^h}{p_t^H}\right)^{\alpha^H}\left(\frac{(1-\alpha^H)\eta_t^h}{e_tp_t^{F*}\!}\right)^{1-\alpha^H} \label{eq:appendix_cpi_derivation_sub4_multiplier_step1} \\
&=\frac{(\eta_t^h)^{\alpha^H+(1-\alpha^H)}}{(\alpha^H)^{\alpha^H}(1-\alpha^H)^{1-\alpha^H}}\cdot\frac{(\alpha^H)^{\alpha^H}(1-\alpha^H)^{1-\alpha^H}}{(p_t^H)^{\alpha^H}(e_tp_t^{F*})^{1-\alpha^H}} \label{eq:appendix_cpi_derivation_sub4_multiplier_step2} \\
c_t^{h\to W}&=\eta_t^h\cdot\frac{1}{(p_t^H)^{\alpha^H}(e_tp_t^{F*})^{1-\alpha^H}} \label{eq:appendix_cpi_derivation_sub4_multiplier_step3}
\end{align}
この式をラグランジュ乗数\(\eta_t^h\)について解く。
\begin{equation}
\eta_t^h=c_t^{h\to W}(p_t^H)^{\alpha^H}(e_tp_t^{F*})^{1-\alpha^H}
\label{eq:appendix_cpi_derivation_sub4_multiplier_result}
\end{equation}

\subsection*{5.総合物価指数\(p_t^{H\to W}\)と最終的な需要関数の導出}
\label{sec:appendix_cpi_derivation_sub5}
名目総消費は、最小化すべき総費用\(p_t^Hc_t^{h\to H}+e_tp_t^{F*}c_t^{h\to F}\)として定義される。ステップ2のFOCより、これは\(\alpha^H\eta_t^h+(1-\alpha^H)\eta_t^h=\eta_t^h\)に等しい。一方で、名目総消費は\(p_t^{H\to W}c_t^{h\to W}\)とも定義されるため、\(p_t^{H\to W}c_t^{h\to W}=\eta_t^h\)という関係が成立しなければならない。

この関係式にステップ4で導出した\(\eta_t^h\)の式を代入する。
\begin{equation}
p_t^{H\to W}c_t^{h\to W}=c_t^{h\to W}(p_t^H)^{\alpha^H}(e_tp_t^{F*})^{1-\alpha^H}
\label{eq:appendix_cpi_derivation_sub5_equivalence}
\end{equation}
両辺の\(c_t^{h\to W}\)を消去することで、定理の項目2で示された\textbf{総合物価指数\(p_t^{H\to W}\)}が導かれる。
\begin{equation}
p_t^{H\to W}=(p_t^H)^{\alpha^H}(e_tp_t^{F*})^{1-\alpha^H}
\label{eq:appendix_cpi_derivation_sub5_cpi_final}
\end{equation}
最後に、この\(\eta_t^h=p_t^{H\to W}c_t^{h\to W}\)という関係をステップ3で得た需要関数に代入することで、定理の項目1で示された最終的な財バスケットへの需要関数が完成する。
\begin{equation}
c_t^{h\to H}=\frac{\alpha^H(p_t^{H\to W}c_t^{h\to W})}{p_t^H}=\alpha^H\frac{p_t^{H\to W}}{p_t^H}c_t^{h\to W}
\label{eq:appendix_cpi_derivation_sub5_demand_domestic_final}
\end{equation}
\begin{equation}
c_t^{h\to F}=\frac{(1-\alpha^H)(p_t^{H\to W}c_t^{h\to W})}{e_tp_t^{F*}}=(1-\alpha^H)\frac{p_t^{H\to W}}{e_tp_t^{F*}}c_t^{h\to W}
\label{eq:appendix_cpi_derivation_sub5_demand_foreign_final}
\end{equation}
（証明終）         % CPI の定義
% !TeX root = ../../main.tex
% sections/app/appendix_production_aggregation.tex

\chapter{生産関数の集計と価格分散}
\label{chap:appendix_production_aggregation}

本付録では、個々の家計の生産関数から出発し、価格の異質性が存在する経済における国全体の集計生産関数を導出する。

\section*{定理5：価格分散を考慮した集計生産関数}
\label{sec:appendix_production_aggregation_theorem5}

個人の生産関数が\(y_t^h=a_t^Hl_t^h\)であり、個別財への需要が\(y_t^h=(p_t^h/p_t^H)^{-\theta^H}Y_t^H\)で与えられる経済において、国全体の集計生産関数は以下のように表される。
\begin{equation}
Y_t^H=\frac{a_t^HL_t^H}{\Delta_t^H}
\label{eq:appendix_production_aggregation_theorem5_main}
\end{equation}
ここで、\(L_t^H\)は総労働投入量\(\sum_{h\in H}l_t^h\)であり、\(\Delta_t^H\)は経済全体の非効率性を示す価格分散項であり、次式で定義される。
\begin{equation}
\Delta_t^H\equiv\sum_{h\in H}\left(\frac{p_t^h}{p_t^H}\right)^{-\theta^H}
\label{eq:appendix_production_aggregation_theorem5_dispersion_def}
\end{equation}

\section*{証明}
\label{sec:appendix_production_aggregation_proof}

この定理を証明するために、総労働投入量\(L_t^H\)の定義式から出発し、集計生産関数を導出する。

\subsection*{1.総労働の定義式からの展開}
\label{sec:appendix_production_aggregation_sub1}
国全体の総労働投入量\(L_t^H\)は、全ての家計の労働投入量の合計である。
\begin{equation}
L_t^H=\sum_{h\in H}l_t^h
\label{eq:appendix_production_aggregation_sub1_labor_def}
\end{equation}
この式に、個人の生産関数\(y_t^h=a_t^Hl_t^h\)を\(l_t^h\)について解いた\(l_t^h=y_t^h/a_t^H\)を代入する。国全体の生産性\(a_t^H\)は全ての家計で共通であるため、総和の外に出すことができる。
\begin{align}
L_t^H&=\sum_{h\in H}\frac{y_t^h}{a_t^H} \label{eq:appendix_production_aggregation_sub1_subst1_step} \\
&=\frac{1}{a_t^H}\sum_{h\in H}y_t^h
\label{eq:appendix_production_aggregation_sub1_subst1_final}
\end{align}
次に、個別財への需要関数\(y_t^h=\left(p_t^h/p_t^H\right)^{-\theta^H}Y_t^H\)を代入する。集計産出量\(Y_t^H\)は個々の家計\(h\)に依存しないため、これも総和の外に出すことができる。
\begin{align}
L_t^H&=\frac{1}{a_t^H}\sum_{h\in H}\left[\left(\frac{p_t^h}{p_t^H}\right)^{-\theta^H}Y_t^H\right] \label{eq:appendix_production_aggregation_sub1_subst2_step} \\
&=\frac{Y_t^H}{a_t^H}\sum_{h\in H}\left(\frac{p_t^h}{p_t^H}\right)^{-\theta^H}
\label{eq:appendix_production_aggregation_sub1_subst2_final}
\end{align}

\subsection*{2.価格分散項の定義と結論}
\label{sec:appendix_production_aggregation_sub2}
ここで、総和の部分を価格分散項\(\Delta_t^H\)として定義する。
\begin{equation}
\Delta_t^H\equiv\sum_{h\in H}\left(\frac{p_t^h}{p_t^H}\right)^{-\theta^H}
\label{eq:appendix_production_aggregation_sub2_dispersion_id}
\end{equation}
この定義をステップ1で得られた式に代入すると、総労働\(L_t^H\)と総生産\(Y_t^H\)の間に以下の厳密な関係式が成立する。
\begin{equation}
L_t^H=\frac{Y_t^H}{a_t^H}\Delta_t^H
\label{eq:appendix_production_aggregation_sub2_relation_pre}
\end{equation}
この式を\(Y_t^H\)について解くことで、定理で示された集計生産関数が得られる（証明終）。
\begin{equation}
Y_t^H=\frac{a_t^HL_t^H}{\Delta_t^H}
\label{eq:appendix_production_aggregation_sub2_final_result}
\end{equation} % 生産関数の集計
%\input{chapters/sections/appendix/aggregate_output.tex}      % 集計生産量と資源制約
% !TeX root = ../../main.tex
% sections/app/appendix_equivalence.tex

\chapter{効用最大化問題と2段階最適化の等価性}
\label{chap:appendix_equivalence}

本付録では、本稿のモデル構築において採用している「個別財の費用最小化問題を解き、そこで得られた価格指数を用いて効用最大化問題を解く」という2段階の最適化アプローチが、全ての個別財の消費量を一度に選択するという単一の効用最大化問題と数学的に完全に等価であることを証明する。

\section*{定理1：最適化問題の等価性}
\label{sec:appendix_equivalence_theorem1}

家計\(h\)が、ある時点\(t\)において所与の総支出額\(E_t^h\)の下で、全ての個別財\(\{c_t^{h\to h'}\}\),\(\{c_t^{h\to f'}\}\)の消費量を直接選択する単一の効用最大化問題は、以下の2段階の最適化問題と等価である。
\begin{enumerate}
    \item\label{itm:appendix_equivalence_theorem1_step1}\textbf{第1段階（費用最小化）}:所与の財バスケット量\(c_t^{h\to H}\),\(c_t^{h\to F}\)を達成するための最小費用と、その際の個別財への需要を求める。この過程で、真の価格指数\(p_t^H\),\(p_t^{F*}\)が導出される。
    \item\label{itm:appendix_equivalence_theorem1_step2}\textbf{第2段階（効用最大化）}:第1段階で導出された価格指数を所与として、予算制約\(p_t^Hc_t^{h\to H}+e_tp_t^{F*}c_t^{h\to F}=E_t^h\)の下で、財バスケット\(c_t^{h\to H}\),\(c_t^{h\to F}\)の最適な組み合わせを選択し、効用を最大化する。
\end{enumerate}

\section*{証明}
\label{sec:appendix_equivalence_proof}

この定理を証明するために、出発点となる厳密な単一の効用最大化問題を、数学的に等価な変形を施すことで、2段階の最適化問題へと帰着させる。

\subsection*{1.出発点：単一の効用最大化問題}
\label{sec:appendix_equivalence_sub1}
厳密な問題設定は、時点\(t\)において、以下の通りである。
\begin{itemize}
    \item\label{itm:appendix_equivalence_sub1_objective}\textbf{最大化対象}:
    \begin{equation}
    \max_{\{c_t^{h\to h'}\},\{c_t^{h\to f'}\}}\ln(c_t^{h\to W})
    \label{eq:appendix_equivalence_sub1_single_max_objective}
    \end{equation}
    ただし、\(c_t^{h\to W}\),\(c_t^{h\to H}\),\(c_t^{h\to F}\)は各個別財消費量の関数である。

    \item\label{itm:appendix_equivalence_sub1_constraint}\textbf{制約条件}:
    \begin{equation}
    \sum_{h'\in H}p_t^{h'}c_t^{h\to h'}+\sum_{f'\in F}p_t^{f'}c_t^{h\to f'}=E_t^h
    \label{eq:appendix_equivalence_sub1_single_max_constraint}
    \end{equation}
\end{itemize}

\subsection*{2.問題の分解}
\label{sec:appendix_equivalence_sub2}
この最大化問題は、総支出\(E_t^h\)を国内財への支出\(E_t^{h\to H}\)と海外財への支出\(E_t^{h\to F}\)に分割する、以下のネストした（入れ子構造の）問題と等価である。
\begin{equation}
\max_{E_t^{h\to H},E_t^{h\to F}}\left(\max_{\{c_t^{h\to h'}\},\{c_t^{h\to f'}\}}\ln(c_t^{h\to W})\quad\text{s.t.}\sum p_t^{h'}c_t^{h\to h'}=E_t^{h\to H},\sum p_t^{f'}c_t^{h\to f'}=E_t^{h\to F}\right)
\label{eq:appendix_equivalence_sub2_nested_problem}
\end{equation}
\begin{equation}
\text{subject to}\quad E_t^{h\to H}+E_t^{h\to F}=E_t^h
\label{eq:appendix_equivalence_sub2_nested_constraint}
\end{equation}

\subsection*{3.双対性（duality）の利用}
\label{sec:appendix_equivalence_sub3}
ここで、内側の最大化問題に注目する。例えば国内財については、「所与の支出\(E_t^{h\to H}\)で、バスケット\(c_t^{h\to H}\)の量を最大化する問題」である。
\begin{equation}
\max_{\{c_t^{h\to h'}\}}c_t^{h\to H}(\{c_t^{h\to h'}\})\quad\text{s.t.}\quad\sum_{h'\in H}p_t^{h'}c_t^{h\to h'}=E_t^{h\to H}
\label{eq:appendix_equivalence_sub3_duality_max}
\end{equation}
ミクロ経済学の双対性（duality）の原理により、この問題は、「所与のバスケット量\(\bar{c}_t^{h\to H}\)を、最小の費用で達成する問題」と完全に等価である。
\begin{equation}
\min_{\{c_t^{h\to h'}\}}\sum_{h'\in H}p_t^{h'}c_t^{h\to h'}\quad\text{s.t.}\quad c_t^{h\to H}(\{c_t^{h\to h'}\})=\bar{c}_t^{h\to H}
\label{eq:appendix_equivalence_sub3_duality_min}
\end{equation}
この費用最小化問題こそが、2段階アプローチにおける第1段階に他ならない。本稿の付録\ref{chap:appendix_cost_minimization}で示したように、この問題を解くことで、所与のバスケット量\(c_t^{h\to H}\)を達成するための最小費用（支出関数）は\(p_t^Hc_t^{h\to H}\)となることが導かれる。同様に、外国財バスケットの最小費用は\(e_tp_t^{F*}c_t^{h\to F}\)となる。

\subsection*{4.問題の再定式化と結論}
\label{sec:appendix_equivalence_sub4}
この支出関数の関係\(E_t^{h\to H}=p_t^Hc_t^{h\to H}\)と\(E_t^{h\to F}=e_tp_t^{F*}c_t^{h\to F}\)を、ステップ2で分解した問題に代入する。すると、問題は以下のように書き換えられる。
\begin{itemize}
    \item\label{itm:appendix_equivalence_sub4_objective}\textbf{最大化}:
    \begin{equation}
    \max_{c_t^{h\to H},c_t^{h\to F}}\ln\left(\frac{(c_t^{h\to H})^{\alpha^H}(c_t^{h\to F})^{1-\alpha^H}}{(\alpha^H)^{\alpha^H}(1-\alpha^H)^{1-\alpha^H}}\right)
    \label{eq:appendix_equivalence_sub4_final_max}
    \end{equation}
    \item\label{itm:appendix_equivalence_sub4_constraint}\textbf{制約条件}:
    \begin{equation}
    p_t^Hc_t^{h\to H}+e_tp_t^{F*}c_t^{h\to F}=E_t^h
    \label{eq:appendix_equivalence_sub4_final_constraint}
    \end{equation}
\end{itemize}
この書き換えられた問題は、まさしく2段階アプローチにおける第2段階そのものである。

以上により、厳密な単一の効用最大化問題が、本稿で採用した2段階の最適化アプローチと数学的に完全に等価であることが証明された（証明終）           % 変数の同値関係（ 再帰的価格等 ）
% !TeX root = ../../main.tex
% sections/app/appendix_resource_constraint.tex

\chapter{国全体の資源制約式の導出}
\label{chap:appendix_resource_constraint}

本付録では、個々の家計の予算制約式を集計し、国内市場の均衡条件を適用することで、国全体の資源制約式を導出する。

\section*{定理7：国全体の資源制約式}
\label{sec:appendix_resource_constraint_theorem7}

国内金融市場が完備であり、政府が均衡予算を達成する経済において、全ての個人の予算制約式を集計すると、
以下の国全体の資源制約式が得られる。
\begin{equation}
\label{form:appendix_resource_constraint_theorem7_main}
p_t^{H\to W}C_t^{H\to W}+B_{t+1}^H=p_t^HY_t^H+(1+i_{t-1}^F)\frac{e_t^{/*}}{e_{t-1}^{/*}}B_t^H
\end{equation}
ここで、\(C_t^{H\to W}\)は国全体の総消費、\(Y_t^H\)は集計産出量、\(B_{t+1}^H\)は期末の対外純資産（自国通貨建て）を表す。

\section*{証明}
\label{sec:appendix_resource_constraint_proof}

\subsection*{1.国全体での集計}
\label{sec:appendix_resource_constraint_sub1}
まず、家計\(h\)の期ごとの名目予算制約式を、
国内の全ての家計\(h\in H\)について足し合わせる（\(\sum_{h\in H}\)）。全ての変数は時点\(t\)に依存する。
\begin{equation}
\label{form:appendix_resource_constraint_sub1_aggregate_sum}
\begin{aligned}
&\sum_{h\in H}\left(\operatorname{E}_t[q_{t,t+1}d_{t+1}^{h\to H}]+b_{t+1}^{h\to F}+p_t^{H\to W}c_t^{h\to W}\right) \\
&\qquad=\sum_{h\in H}\left(d_t^{h\to H}+(1+i_{t-1}^F)\frac{e_t^{/*}}{e_{t-1}^{/*}}b_t^{h\to F}+(1-\tau_t^H)p_t^hy_t^h+t_t^H\right)
\end{aligned}
\end{equation}

\subsection*{2.国内取引の相殺}
\label{sec:appendix_resource_constraint_sub2}
次に、国全体で集計するとゼロになる国内完結の取引を相殺する。

\paragraph{国内債券市場の均衡}
\label{sec:appendix_resource_constraint_sub2_bond}
国内で取引される状態コンティンジェント債券は、国内の誰かの負債が他の誰かの資産となるゼロサム取引であるため、純供給はゼロである。
\begin{equation}
\label{form:appendix_resource_constraint_sub2_bond_clearing}
\sum_{h\in H}d_t^{h\to H}=0\quad\text{and}\quad\sum_{h\in H}\operatorname{E}_t[q_{t,t+1}d_{t+1}^{h\to H}]=0
\end{equation}
これにより、集計した式から国内債券に関する項（\(d\)）はすべて消去される。

\paragraph{政府部門の予算制約}
\label{sec:appendix_resource_constraint_sub2_gov}
政府は均衡予算を達成し、税収のすべてを家計への移転に使うため、移転の総額と税収の総額は一致する。
\begin{equation}
\label{form:appendix_resource_constraint_sub2_gov_budget}
\sum_{h\in H}t_t^H=\sum_{h\in H}\tau_t^Hp_t^hy_t^h
\end{equation}
この関係を用いると、集計後の予算制約式の右辺にある移転項\(\sum t_t^H\)は、所得項の一部である税金部分\(\sum\tau_t^Hp_t^hy_t^h\)と相殺される。

\subsection*{3.集計変数への書き換え}
\label{sec:appendix_resource_constraint_sub3}
国内取引が相殺された結果、式に残るのは以下の項のみである。
\begin{equation}
\label{form:appendix_resource_constraint_sub3_remaining_terms}
\sum_{h\in H}\left(b_{t+1}^{h\to F}+p_t^{H\to W}c_t^{h\to W}\right)=\sum_{h\in H}\left((1+i_{t-1}^F)\frac{e_t^{/*}}{e_{t-1}^{/*}}b_t^{h\to F}+p_t^hy_t^h\right)
\end{equation}
この式を、国全体の集計変数（大文字の変数）を使って書き換える。

\paragraph{消費と対外資産の集計}
\label{sec:appendix_resource_constraint_sub3_agg}
価格指数\(p_t^{H\to W}\)は全ての家計で共通であるため、総和\(\sum_{h\in H}\)の外に出すことができる。
\begin{equation}
\label{form:appendix_resource_constraint_sub3_consumption_agg}
\sum_{h\in H}p_t^{H\to W}c_t^{h\to W}=p_t^{H\to W}\sum_{h\in H}c_t^{h\to W}\equiv p_t^{H\to W}C_t^{H\to W}
\end{equation}
対外純資産（自国通貨建て）については、単純な総和として定義される。
\begin{equation}
\label{form:appendix_resource_constraint_sub3_nfa_def}
\sum_{h\in H}b_{t+1}^{h\to F}\equiv B_{t+1}^H
\end{equation}
同様に、前期からの対外資産の償還額も\((1+i_{t-1}^F)\frac{e_t^{/*}}{e_{t-1}^{/*}}B_t^H\)となる。

\paragraph{名目総生産の集計}
\label{sec:appendix_resource_constraint_sub3_gdp}
付録\ref{chap:appendix_aggregate_output}で証明した通り、総名目所得は\(\sum_{h\in H}p_t^hy_t^h=p_t^HY_t^H\)という関係が厳密に成立する。

\subsection*{4.結論}
\label{sec:appendix_resource_constraint_sub4}
以上の集計結果をまとめると、定理で示された国全体の資源制約式が導出される（証明終）。
\begin{equation}
\label{form:appendix_resource_constraint_sub4_final}
p_t^{H\to W}C_t^{H\to W}+B_{t+1}^H=p_t^HY_t^H+(1+i_{t-1}^F)\frac{e_t^{/*}}{e_{t-1}^{/*}}B_t^H
\end{equation}    % 資源制約の線形化
% !TeX root = ../../main.tex
% sections/app/appendix_risk_sharing.tex

\chapter{完備市場における消費の共通化と貯蓄の個別化}
\label{chap:appendix_risk_sharing}

本付録では、国内完備市場とカルボ型価格設定を仮定したモデルにおいて、なぜ全ての家計の消費が共通化される一方、貯蓄（資産ポートフォリオ）は家計ごとに個別化されるのかを厳密に証明する。

\section*{定理9：リスク共有の結果}
\label{sec:appendix_risk_sharing_theorem9}

国内金融市場が完備である経済において、以下の関係が成立する。
\begin{enumerate}
    \item\label{itm:appendix_risk_sharing_theorem9_lambda}全ての国内家計\(h\in H\)の所得の限界効用は、いかなる時点\(t\)、いかなる状態\(j\)においても常に一致する。
    \begin{equation}
    \lambda_t^h=\lambda_t^{h'}\quad\forall h,h'\in H
    \label{eq:appendix_risk_sharing_theorem9_lambda_equality}
    \end{equation}
    \item\label{itm:appendix_risk_sharing_theorem9_consumption}全ての国内家計\(h\in H\)の総消費指数は、常に一致する。
    \begin{equation}
    c_t^{h\to W}=c_t^{h'\to W}\quad\forall h,h'\in H
    \label{eq:appendix_risk_sharing_theorem9_consumption_equality}
    \end{equation}
    \item\label{itm:appendix_risk_sharing_theorem9_portfolio}各家計が購入する次期のための資産ポートフォリオの総価値は、家計が当期に価格改定の機会を得たか否かによって異なるため、一般に一致しない。
\end{enumerate}

\section*{証明}
\label{sec:appendix_risk_sharing_proof}

\subsection*{1.所得の限界効用\(\lambda_t\)の一致（定理9-1の証明）}
\label{sec:appendix_risk_sharing_sub1}

\paragraph{ステップA：限界効用の比率の不変性}
\label{sec:appendix_risk_sharing_sub1_step_a}
家計の最適化行動は全ての時点・状態で成立するため、任意の2つの国内家計\(h\)と\(h'\)の国内コンティンジェント債券に関する一階の条件（FOC）は、常に成立する。
\begin{align}
\lambda_t^hq_{t+1}&=\beta_t^H\pi\lambda_{t+1}^h \label{eq:appendix_risk_sharing_sub1_foc_h} \\
\lambda_t^{h'}q_{t+1}&=\beta_t^H\pi\lambda_{t+1}^{h'} \label{eq:appendix_risk_sharing_sub1_foc_h_prime}
\end{align}
これら2式の比を取ると共通項が消去され、以下の関係が得られる。
\begin{equation}
\frac{\lambda_t^h}{\lambda_t^{h'}}=\frac{\lambda_{t+1}^h}{\lambda_{t+1}^{h'}}=k
    \label{eq:appendix_risk_sharing_sub1_constant_ratio}
\end{equation}
この比率\(k\)は、時間や状態に依存しない普遍的な定数である。

\paragraph{ステップB：定数\(k\)の値の特定}
\label{sec:appendix_risk_sharing_sub1_step_b}
第\ref{sec:model_overview}節で定義した「事的事実対称性」の仮定を用いる。初期時点\(s\)において、全ての家計は同一の選好をもち、かつ初期貯蓄がゼロ（\(d_s^h=d_s^{h'}=0\)）である。このとき、各家計が直面する最適化問題は数学的に完全に同一であるため、初期の所得の限界効用は全ての家計で一致しなければならない。
\begin{equation}
\lambda_s^h=\lambda_s^{h'}\quad\Longrightarrow\quad k=\frac{\lambda_s^h}{\lambda_s^{h'}}=1
\label{eq:appendix_risk_sharing_sub1_k_identity}
\end{equation}

\paragraph{ステップC：結論}
\label{sec:appendix_risk_sharing_sub1_step_c}
普遍的な定数が\(k=1\)であることから、将来のすべての時点・すべての状態において\(\lambda_t^h=\lambda_t^{h'}\)が成立する。

\subsection*{2.総消費指数\(c_t^{h\to W}\)の一致（定理9-2の証明）}
\label{sec:appendix_risk_sharing_sub2}
全ての家計の\(\lambda_t\)が一致するという結果を、消費に関するFOC（式\ref{eq:model_foc_consumption_home}）に適用する。
\begin{equation}
\lambda_t=\frac{1}{p_t^{H\to W}c_t^{h\to W}}
\label{eq:appendix_risk_sharing_sub2_lambda_def}
\end{equation}
左辺の\(\lambda_t\)と右辺の物価指数\(p_t^{H\to W}\)は全ての家計で共通であるため、総消費指数\(c_t^{h\to W}\)もまた、全ての家計間で完全に一致しなければならない。

\subsection*{3.ポートフォリオ購入総額の個別化（定理9-3の証明）}
\label{sec:appendix_risk_sharing_sub3}
家計\(h\)の予算制約式を、次期のために購入する資産ポートフォリオの総価値\(v_{t+1}^h\equiv\sum_{j'}q_{t+1}d_{t+1}^{h\to H}+e_tb_{t+1}^{h\to F*}\)について整理する。



\begin{equation}
v_{t+1}^h=\underbrace{\left(d_t^{h\to H}+(1+i_{t-1}^F)e_tb_t^{h\to F*}+(1-\tau_t^H)p_t^hy_t^h+t_t^H\right)}_{\text{当期の総収入}}-\underbrace{p_t^{H\to W}c_t^{h\to W}}_{\text{当期の消費支出}}
\label{eq:appendix_risk_sharing_sub3_portfolio_dynamics}
\end{equation}
定理9-2より消費支出は共通であるが、当期の総収入、特に労働所得\((1-\tau_t^H)p_t^hy_t^h\)はカルボ型の価格設定（\(p_t^h\)の異質性）によって家計ごとに異なる。支出が共通で収入が異なる以上、その差額を埋める資産保有額\(v_{t+1}^h\)は、各家計が直面したショックに応じて個別化される。（証明終）          % リスクシェアリング条件

% =============================================================
% 2. 第 4 章（ モデル構築 ）で参照される基礎的な導出
% =============================================================
\input{chapters/sections/appendix/optimal_price_derivation.tex}          % フィリップス曲線

% =============================================================
% 3. 第 5 章（ シミュレーション結果 ）で参照される詳細な分析
% =============================================================
% !TeX root = ../../main.tex
% sections/app/appendix_exchange_rate_and_gdp.tex

\chapter{名目為替レートおよび名目GDPの式の導出}
\label{chap:appendix_exchange_rate_and_gdp}

本付録では名目為替レートおよび名目GDPの式を導出する。
具体的には、自国と外国の代表的家計の
国際債券に関するオイラー方程式、自国および外国債券に関するオイラー方程式、
所得の限界効用の方程式より為替レートの式を特定し、
それを財市場の均衡条件と組み合わせることで名目GDPの式を導出する。
さらにこれらの関係式を用いることにより、
全期間において自国と外国の主観的割引因子が等しい（\(\beta_t^H=\beta_t^F\)）場合には、
第\ref{sec:model_steady_state}節で仮定した初期条件\(b_s^H=0\)のもとで、
帰納的にすべての\(t\ge s\)において対外純資産が恒等的にゼロ（\(b_t^H=0\)）となることを証明する。

\section*{1.期待値の比に関する線形近似の証明}
\label{sec:appendix_exchange_rate_and_gdp_s1}

本稿のような線形近似モデルにおいて、
期待値の比\(E_t[X_t]/E_t[Y_t]\)が比の期待値\(E_t[X_t/Y_t]\)で近似できる理由を
一般的な確率変数の性質にもとづいて数学的に証明する。

\subsection*{一般論としての変数の定義}
\label{sec:appendix_exchange_rate_and_gdp_s1_def}
ある任意の確率変数\(X_t\)について、その期待値（平均）を\(\bar{X}\equiv E[X_t]\)とし、期待値からの対数乖離（変化率）を\(\hat{x}_t\)と定義する。
\begin{equation}
    \hat{x}_t\equiv\frac{X_t-\bar{X}}{\bar{X}}\quad\Longleftrightarrow\quad X_t=\bar{X}(1+\hat{x}_t)
    \label{eq:appendix_exchange_rate_and_gdp_s1_log_linear_def}
\end{equation}
本稿のような線形近似モデルでは、この乖離\(\hat{x}_t\)を「1次の微小量」として扱い、その2乗以上（\(\hat{x}_t^2\)や\(\hat{x}_t\hat{y}_t\)）を無視できるほど小さい項（\(\approx0\)）として計算を行う。

\subsection*{ステップ1：積の期待値の近似（共分散項の評価）}
\label{sec:appendix_exchange_rate_and_gdp_s1_sub1}

一般に、2つの確率変数\(X,Y\)の積の期待値は「期待値の積」と「共分散」の和に分解できる。
\begin{equation}
    E_t[XY]=E_t[X]E_t[Y]+\text{Cov}_t(X,Y)
    \label{eq:appendix_exchange_rate_and_gdp_s1_covariance_decomposition}
\end{equation}
ここで、共分散の項を上記の定義に従って展開する。
\begin{align}
    \text{Cov}_t(X,Y)&=E_t\left[(X-\bar{X})(Y-\bar{Y})\right] \nonumber \\
    &=E_t\left[\left(\bar{X}(1+\hat{x})-\bar{X}\right)\left(\bar{Y}(1+\hat{y})-\bar{Y}\right)\right] \nonumber \\
    &=\bar{X}\bar{Y}E_t\left[\hat{x}\hat{y}\right]
    \label{eq:appendix_exchange_rate_and_gdp_s1_covariance_order}
\end{align}
ここで、右辺の期待値の中身は微小変動同士の積（\(\hat{x}\times\hat{y}\)のオーダー）となっている。1次近似の枠組みでは、1次の微小量同士の積は2次の微小量となるため、無視することができる（\(\approx0\)）。したがって、以下の近似が成立する。
\begin{equation}
    \text{Cov}_t(X,Y)\approx0\quad\Longrightarrow\quad E_t[XY]\approx E_t[X]E_t[Y]
    \label{eq:appendix_exchange_rate_and_gdp_s1_approx_product}
\end{equation}

\subsection*{ステップ2：逆数の期待値の近似（分散項の評価）}
\label{sec:appendix_exchange_rate_and_gdp_s1_sub2}

次に、変数\(Z\)の逆数の期待値\(E_t[1/Z]\)を考える。一般に、イェンゼンの不等式により\(E[1/Z]\neq1/E[Z]\)であるが、その乖離幅は分散に依存する。
関数\(f(Z)=1/Z\)を期待値\(\bar{Z}=E_t[Z]\)の周りで2次までテイラー展開し期待値をとる。
\begin{align}
    E_t\left[\frac{1}{Z}\right]&\approx E_t\left[\frac{1}{\bar{Z}}-\frac{1}{\bar{Z}^2}(Z-\bar{Z})+\frac{1}{\bar{Z}^3}(Z-\bar{Z})^2\right] \nonumber \\
    &=\frac{1}{\bar{Z}}-\frac{1}{\bar{Z}^2}\underbrace{E_t[Z-\bar{Z}]}_{0}+\frac{1}{\bar{Z}^3}\underbrace{E_t[(Z-\bar{Z})^2]}_{\text{分散Var}(Z)} \nonumber \\
    &=\frac{1}{E_t[Z]}+\frac{1}{(E_t[Z])^3}\text{Var}_t(Z)
    \label{eq:appendix_exchange_rate_and_gdp_s1_jensen_approx}
\end{align}
ここで分散\(\text{Var}_t(Z)=\bar{Z}^2E_t\left[\hat{z}^2\right]\)は微小変動の2乗を含むため、1次近似においてはこれを無視できる。
したがって以下の近似式が成立する。
\begin{equation}
    E_t\left[\frac{1}{Z}\right]\approx\frac{1}{E_t[Z]}
    \label{eq:appendix_exchange_rate_and_gdp_s1_approx_inverse}
\end{equation}

\subsection*{結論：比の期待値への変換}
\label{sec:appendix_exchange_rate_and_gdp_s1_conclusion}

式\eqref{eq:appendix_exchange_rate_and_gdp_s1_approx_product}と式\eqref{eq:appendix_exchange_rate_and_gdp_s1_approx_inverse}の結果を組み合わせることで、期待値の比を比の期待値として扱うことが正当化される。
\begin{align}
    \frac{E_t[X]}{E_t[Y]}&=E_t[X]\cdot\frac{1}{E_t[Y]} \nonumber \\
    &\approx E_t[X]\cdot E_t\left[\frac{1}{Y}\right]\quad（\because\text{分散項}\approx0） \nonumber \\
    &\approx E_t\left[X\cdot\frac{1}{Y}\right]\quad（\because\text{共分散項}\approx0） \nonumber \\
    &=E_t\left[\frac{X}{Y}\right]
    \label{eq:appendix_exchange_rate_and_gdp_s1_ratio_final}
\end{align}
本付録の以降の導出では、この一般的な近似式\(E_t[X_t]/E_t[Y_t]\approx E_t[X_t/Y_t]\)を利用する。

\section*{2.自国と外国の代表的家計の国際債券に関する修正オイラー方程式の比}
\label{sec:appendix_exchange_rate_and_gdp_s2}
第\ref{sec:appendix_insulation_effect_modified_uip}節で確認したとおり、
自国と外国の代表的家計の国際債券に関する修正オイラー方程式は次のようであった。
\begin{equation}
    \lambda_t^{H/*}=\beta_t^H(1+i_t^F)E_t[\lambda_{t+1}^{H/*}]
    \tag{\ref{eq:appendix_insulation_effect_modified_final_uip_H}}
\end{equation}
\begin{equation}
    \lambda_t^{F/*}=\beta_t^F(1+i_t^F)E_t[\lambda_{t+1}^{F/*}]
    \tag{\ref{eq:appendix_insulation_effect_modified_final_uip_F}}
\end{equation}

ここで自国と外国の代表的家計の外国通貨建て所得の限界効用の比率を\(\Lambda_t\)と定義する。
\begin{equation}
    \Lambda_t\equiv\frac{\lambda_t^{H/*}}{\lambda_t^{F/* messenger}}
    \label{eq:appendix_exchange_rate_and_gdp_s2_lambda_ratio_def}
\end{equation}
式\eqref{eq:appendix_insulation_effect_modified_final_uip_H}と
式\eqref{eq:appendix_insulation_effect_modified_final_uip_F}
の辺々を割ると、共通項である外国の名目利子率\((1+i_t^F)\)が相殺され以下の関係が得られる。
\begin{equation}
    \Lambda_t=\frac{\beta_t^H}{\beta_t^F}\frac{E_t[\lambda_{t+1}^{H/*}]}{E_t[\lambda_{t+1}^{F/*}]}
    \label{eq:appendix_exchange_rate_and_gdp_s2_lambda_ratio_raw}
\end{equation}

第1節で証明した一般的な近似関係を用いることで、式\eqref{eq:appendix_exchange_rate_and_gdp_s2_lambda_ratio_raw}の右辺にある「期待値の比」を「比の期待値」へと変換し、以下の再帰式を得る。
\begin{equation}
    \Lambda_t\approx\frac{\beta_t^H}{\beta_t^F}E_t\left[\frac{\lambda_{t+1}^{H/*}}{\lambda_{t+1}^{F/*}}\right]=\frac{\beta_t^H}{\beta_t^F}E_t[\Lambda_{t+1}]
    \label{eq:appendix_exchange_rate_and_gdp_s2_lambda_recursive}
\end{equation}

\section*{3.将来に向けた前方展開と極限}
\label{sec:appendix_exchange_rate_and_gdp_s3}

前節で導出した再帰式\eqref{eq:appendix_exchange_rate_and_gdp_s2_lambda_recursive}を将来に向かって逐次代入することで、現在の比率\(\Lambda_t\)を決定する。

\paragraph{①再帰的代入と一般形の導出}
\label{sec:appendix_exchange_rate_and_gdp_s3_p1}
式\eqref{eq:appendix_exchange_rate_and_gdp_s2_lambda_recursive}の右辺にある\(\Lambda_{t+1}\)に対して、1期先の関係式\(\Lambda_{t+1}=\frac{\beta_{t+1}^H}{\beta_{t+1}^F}E_{1+t}[\Lambda_{t+2}]\)を代入する。この操作を\(k\)回繰り返すと、反復期待値の法則により、以下の一般形が得られる。
\begin{equation}
    \Lambda_t=E_t\left[\left(\prod_{j=0}^{k}\frac{\beta_{t+j}^H}{\beta_{t+j}^F}\right)\Lambda_{t+k+1}\right]
    \label{eq:appendix_exchange_rate_and_gdp_s3_iterative_general}
\end{equation}

\paragraph{②極限の適用}
\label{sec:appendix_exchange_rate_and_gdp_s3_p2}
次に、式\eqref{eq:appendix_exchange_rate_and_gdp_s3_iterative_general}の両辺について\(k\to\infty\)の極限をとる。
本稿のモデルにおいては、全ての変数が定常状態近傍で安定的に推移（有界性）することを前提としている。
このときルベーグの優収束定理により
期待値演算子\(E_t\)と極限操作\(\lim\)の順序を入れ替えることが正当化される。
すなわち、期待値の極限は極限の期待値に等しい。

\begin{align}
    \Lambda_t&=\lim_{k\to\infty}E_t\left[\left(\prod_{j=0}^{k}\frac{\beta_{t+j}^H}{\beta_{t+j}^F}\right)\Lambda_{t+k+1}\right] \nonumber \\
    &=E_t\left[\lim_{k\to\infty}\left(\left(\prod_{j=0}^{k}\frac{\beta_{t+j}^H}{\beta_{t+j}^F}\right)\Lambda_{t+k+1}\right)\right]
    \label{eq:appendix_exchange_rate_and_gdp_s3_limit_exchange}
\end{align}

ここで、経済の安定性条件より、無限遠方において一時的なショックの影響は消失し、
変数はショック後の新たな定常状態へ収束すると仮定する。
すると\(\Lambda_{t+k+1}\)は定常状態における\(\Lambda_{\infty}\)へと収束する。
\begin{equation}
    \lim_{k\to\infty}\Lambda_{t+k+1}=\Lambda_{\infty}
    \label{eq:appendix_exchange_rate_and_gdp_s3_lambda_convergence}
\end{equation}
不完備市場モデルにおいてはこの\(\Lambda_{\infty}\)は
ショックの履歴に依存する確率変数としての性質を持つため、期待値演算子の中に留まる。

ここで資本移動要因\(\mathcal{B}_t\)を以下のように定義する。
\begin{equation}
    \mathcal{B}_t\equiv\prod_{j=0}^{\infty}\frac{\beta_{t+j}^H}{\beta_{t+j}^F}
    \label{eq:appendix_exchange_rate_and_gdp_s3_bt_def}
\end{equation}

以上より\(\Lambda_t\)は以下の形式で確定する。
\begin{equation}
    \Lambda_t=E_t\left[\mathcal{B}_t\Lambda_{\infty}\right]
    \label{eq:appendix_exchange_rate_and_gdp_s3_lambda_final_solution}
\end{equation}

\section*{4.結論：為替レートの決定式}
\label{sec:appendix_exchange_rate_and_gdp_s4}

\(\lambda_t^{H/*}\)の定義式\eqref{eq:appendix_exchange_rate_lambda_def}を\(e_t^{/*}\)について解いた
\(e_t^{/*}=\lambda_t^{H/*}/\lambda_t^H\)に
\(\Lambda_t\)の定義式\eqref{eq:appendix_exchange_lambda_ratio_def}を変形した
\(\lambda_t^{H/*}=\Lambda_t\lambda_t^{F/*}\)を代入すると以下が得られる。
\begin{equation}
    e_t^{/*}=\Lambda_t\times\frac{\lambda_t^{F/*}}{\lambda_t^H}
    \label{eq:appendix_exchange_rate_and_gdp_s4_mu_ratio}
\end{equation}

ここで式\eqref{eq:appendix_exchange_rate_and_gdp_s4_mu_ratio}に
代表的家計のオイラー方程式\eqref{eq:final_euler_H}、\eqref{eq:final_euler_F}を代入すると、
以下の名目利子率を含む関係式が得られる。
\begin{equation}
    e_t^{/*}=\Lambda_t\times\frac{\beta_t^{F*}(1+i_t^F)E_t[\lambda_{t+1}^{F/*}]}{\beta_t^H(1+i_t^H)E_t[\lambda_{t+1}^H]}
    \label{eq:appendix_exchange_rate_and_gdp_s4_euler_ratio}
\end{equation}

一方で式\eqref{eq:appendix_exchange_rate_and_gdp_s4_mu_ratio}に
所得の限界効用の式\eqref{eq:final_lambda_H}、\eqref{eq:final_lambda_F}を代入して整理すると、
最終的な為替レート決定式が得られる。
\begin{align}
    e_t^{/*}&=\Lambda_t\times\frac{\lambda_t^{F/*}}{\lambda_t^H} \nonumber \\
    &=\Lambda_t\times\frac{1/(p_t^{F\to W*}c_t^{F\to W})}{1/(p_t^{H\to W}c_t^{H\to W})} \nonumber \\
    &=E_t[\mathcal{B}_t\Lambda_{\infty}]\times\frac{p_t^{H\to W}c_t^{H\to W}}{p_t^{F\to W*}c_t^{F\to W}}
    \label{eq:appendix_exchange_rate_and_gdp_s4_rate_final}
\end{align}

ここで、期待値項\(E_t\left[\mathcal{B}_t\Lambda_{\infty}\right]\)について、定常状態における割引因子の設定が決定的な役割を果たす。

\begin{itemize}
    \item \textbf{定常状態で割引因子が等しい場合（\(\beta_{ss}^H=\beta_{ss}^F\)）：}
    本稿の基本設定である。一時的なショックにより\(\beta_t\)が変動しても、長期的には元の水準に戻るため、無限乗積項\(\mathcal{B}_t\)は有限の値に収束する。
    特に、全期間において\(\beta_t^H=\beta_t^F\)であるならば、式\eqref{eq:appendix_exchange_rate_and_gdp_s4_rate_final}の期待値項は\(E_t[\Lambda_{\infty}]\)となり、為替レートは純粋に両国の名目支出比率のみで決定される。
    \begin{equation}
        e_t^{/*}=E_t[\Lambda_{\infty}]\times\frac{p_t^{H\to W}c_t^{H\to W}}{p_t^{F\to W*}c_t^{F\to W}}
        \label{eq:appendix_exchange_rate_and_gdp_s4_rate_simplified}
    \end{equation}
    
    \item \textbf{定常状態で割引因子が異なる場合（\(\beta_{ss}^H\neq\beta_{ss}^F\)）：}
    もし恒久的に\(\beta\)が異なると仮定すると、無限乗積項は発散または消失し、安定的な均衡が存在しなくなる。したがって、本モデルのような無限期間モデルにおいて安定解を得るためには、定常状態において両国の割引因子が一致していることが前提条件となる。
\end{itemize}

\section*{5.資源制約式と対外純資産の恒等的なゼロ均衡}
\label{sec:appendix_exchange_rate_and_gdp_s5}

上記の為替レート決定メカニズムが、対外純資産（\(b_t^H\)）の動学に与える影響を確認する。

\subsection*{1.名目GDP方程式の導出と為替レート決定式の代入}
\label{sec:appendix_exchange_rate_and_gdp_s5_sub1}

自国の一人当たり財市場均衡条件\eqref{eq:model_goods_market_eq_H}から出発する。
\begin{equation}
    y_t^H=c_t^{H\to H}+\frac{M}{N}c_t^{F\to H}
    \label{eq:appendix_exchange_rate_and_gdp_s5_sub1_goods_raw}
\end{equation}
この両辺に自国財価格\(p_t^H\)を乗じると、名目GDPは各主体の需要内訳として以下のように展開される。
\begin{align}
    p_t^Hy_t^H&=p_t^Hc_t^{H\to H}+p_t^H\frac{M}{N}c_t^{F\to H} \nonumber \\
    &=p_t^H\left(\alpha^H\frac{p_t^{H\to W}}{p_t^H}c_t^{H\to W}\right)+p_t^H\frac{M}{N}\left((1-\alpha^F)\frac{p_t^{F\to W*}}{p_t^{H*}}c_t^{F\to W}\right) \nonumber \\
    &=\alpha^H(p_t^{H\to W}c_t^{H\to W})+\frac{M}{N}(1-\alpha^F)\frac{p_t^H}{p_t^{H*}}(p_t^{F\to W*}c_t^{F\to W})
    \label{eq:appendix_exchange_rate_and_gdp_s5_sub1_gdp_derivation}
\end{align}
ここで、一物一価の法則\(p_t^H=e_t^{/*}p_t^{H*}\)\eqref{eq:model_lop}および、人口比\(N=M=1\)の設定を適用すると、以下の一般的な名目GDPの方程式が得られる。
\begin{equation}
    p_t^Hy_t^H=\alpha^H(p_t^{H\to W}c_t^{H\to W})+(1-\alpha^F)e_t^{/*}(p_t^{F\to W*}c_t^{F\to W})
    \label{eq:appendix_exchange_rate_and_gdp_s5_sub1_gdp_general}
\end{equation}

さらに、式\eqref{eq:appendix_exchange_rate_and_gdp_s5_sub1_gdp_general}の右辺にある為替レート\(e_t^{/*}\)に、先に導出した決定式\eqref{eq:appendix_exchange_rate_and_gdp_s4_rate_final}を代入する。
\begin{align}
    p_t^Hy_t^H&=\alpha^H(p_t^{H\to W}c_t^{H\to W})+(1-\alpha^F)\left(E_t\left[\mathcal{B}_t\Lambda_{\infty}\right]\frac{p_t^{H\to W}c_t^{H\to W}}{p_t^{F\to W*}c_t^{F\to W}}\right)(p_t^{F\to W*}c_t^{F\to W}) \nonumber \\
    &=\alpha^H(p_t^{H\to W}c_t^{H\to W})+(1-\alpha^F)E_t\left[\mathcal{B}_t\Lambda_{\infty}\right](p_t^{H\to W}c_t^{H\to W}) \nonumber \\
    &=\left[\alpha^H+(1-\alpha^F)E_t\left[\mathcal{B}_t\Lambda_{\infty}\right]\right]p_t^{H\to W}c_t^{H\to W}
    \label{eq:appendix_exchange_rate_and_gdp_s5_sub1_gdp_substituted}
\end{align}

ここで第\ref{sec:model_steady_state}節において仮定されたとおり
初期時点およびそれ以前の対外純資産はゼロ（\(b_s^H=b_{s-1}^H=0\)）である。

自国の資源制約式\eqref{eq:model_resource_constraint_H}においてショック発生前の\(t=s-1\)の時点を考える。
\begin{equation}
    p_{s-1}^{H\to W}c_{s-1}^{H\to W}+b_s^H=p_{s-1}^Hy_{s-1}^H+(1+i_{s-2}^F)\frac{e_{s-1}^{/*}}{e_{s-2}^{/*}}b_{s-1}^H
    \label{eq:appendix_exchange_rate_and_gdp_s5_sub1_res_init}
\end{equation}
ここに仮定\(b_s^H=b_{s-1}^H=0\)を代入すると、以下の貿易収支の均衡条件が導かれる。
\begin{equation}
    p_{s-1}^Hy_{s-1}^H=p_{s-1}^{H\to W}c_{s-1}^{H\to W}
    \label{eq:appendix_exchange_rate_and_gdp_s5_sub1_trade_init}
\end{equation}
式\eqref{eq:appendix_exchange_rate_and_gdp_s5_sub1_gdp_substituted}を初期時点\(t=s-1\)で評価し、貿易収支均衡\eqref{eq:appendix_exchange_rate_and_gdp_s5_sub1_trade_init}と比較すると、右辺の名目消費にかかる係数部分は1でなければならないことがわかる。すなわち、
\begin{equation}
    1=\alpha^H+(1-\alpha^F)E_{s-1}\left[\mathcal{B}_{s-1}\Lambda_{\infty}\right]
    \label{eq:appendix_exchange_rate_and_gdp_s5_sub1_unity}
\end{equation}
が成立する。

\subsection*{2.aショックが発生したときの均衡}
\label{sec:appendix_exchange_rate_and_gdp_s5_sub2}

次に、全期間において自国と外国の割引因子が等しく（\(\beta_t^H=\beta_t^F\)）、生産性ショック（aショック）のみが発生するケースを検討する。

まず、本稿の第5章「生産性ショックにおける国際的波及（遮断効果）」で論じたように、外貨建ての限界効用\(\lambda_t^{H/*}\)および\(\lambda_t^{F/*}\)は生産性ショックの影響を受けない方程式系によって決定される。したがって、ショックの前後でこれらの値は変化せず、その比率である\(\Lambda_t=\lambda_t^{H/*}/\lambda_t^{F/*}\)も不変に維持される。すなわち、
\begin{equation}
    \Lambda_{s-1}=\Lambda_s=\dots=\Lambda_{\infty}
    \label{eq:appendix_exchange_rate_and_gdp_s5_sub2_lambda_constancy}
\end{equation}
が成立する。この性質により、評価の比率の極限値\(\Lambda_{\infty}\)は、ショックの履歴に依存しない確定的定数とみなすことができ、期待値演算子\(E_t\)の外に出すことが可能となる。

また、本ケースでは全期間において割引因子が一定であるため、定義式\eqref{eq:appendix_exchange_rate_and_gdp_s3_bt_def}より資本移動要因は恒等的に\(\mathcal{B}_t=1\)となる。

この条件を式\eqref{eq:appendix_exchange_rate_and_gdp_s5_sub1_unity}に適用すると、
\begin{equation}
    1=\alpha^H+(1-\alpha^F)\Lambda_{\infty}
    \label{eq:appendix_exchange_rate_and_gdp_s5_sub2_unity_subst}
\end{equation}
となり、これを整理すると以下が得られる。
\begin{equation}
    \Lambda_{\infty}=\frac{1-\alpha^H}{1-\alpha^F}
    \label{eq:appendix_exchange_rate_and_gdp_s5_sub2_lambda_ss}
\end{equation}

この\(\Lambda_{\infty}\)と\(\mathcal{B}_t=1\)を式\eqref{eq:appendix_exchange_rate_and_gdp_s5_sub1_gdp_substituted}に代入すると、以下の名目貿易収支の均衡式が得られる。
\begin{equation}
    p_t^Hy_t^H=p_t^{H\to W}c_t^{H\to W}
    \label{eq:appendix_exchange_rate_and_gdp_s5_sub2_trade_balance}
\end{equation}

最後に、式\eqref{eq:appendix_exchange_rate_and_gdp_s5_sub2_trade_balance}を自国の資源制約式\eqref{eq:model_resource_constraint_H}に代入すると、資産蓄積に関する以下の差分方程式が得られる。
\begin{equation}
    b_{t+1}^H=(1+i_{t-1}^F)\frac{e_t^{/*}}{e_{t-1}^{/*}}b_t^H
    \label{eq:appendix_exchange_rate_and_gdp_s5_sub2_nfa_diff}
\end{equation}

初期条件\(b_s^H=0\)を用いると、帰納的にすべての\(t\ge s\)において\(b_{t+1}^H=0\)であることが証明される。以上により、Cole-Obstfeld条件下では為替レートの調整により貿易収支が常に均衡し、対外純資産が恒等的にゼロとなることが数学的に裏付けられた。          % 為替レートと GDP の恒等式
\input{chapters/sections/appendix/log_linearization.tex}      % 線形化フィリップス曲線の導出
% !TeX root = ../../main.tex
% sections/app/appendix_welfare_correction.tex

\chapter{厚生評価における近似の整合性と補正項の導出}
\label{chap:appendix_welfare_correction}

本付録では、本稿のシミュレーション分析で採用している「1次近似（対数線形近似）による動学算出」と「2次近似得的視点に基づく厚生評価」の間の論理的整合性、および具体的な補正ロジックについて詳述する。

\section{近似による情報の欠落と整合性の問題}
\label{sec:appendix_welfare_correction_consistency}

本稿のシミュレーションは対数線形近似（1次近似）を用いているが、この手法で得られた結果を厚生評価に用いる際には、「非線形な関数」と「近似された変数」の取り扱いに細心の注意が必要である。なぜなら、\textbf{「非線形な関数」に「1次近似された変数」をそのまま代入して計算すると、計算の整合性が取れなくなる（偽の精度が生じる）}からである。

この問題を理解するために、簡単な「A+B」の例を考えてみよう。
ある真の値\(X\)が、主要な変動成分である「1次の項\(A\)」と、微細な補正成分である「2次の項\(B\)」から成るとする（\(X=A+B\)）。
一方、評価したい関数が\(F(X)=X+X^2\)という非線形（2次）の関数だとする。

真の値をこの関数に入力して、2次の精度まで正しく計算すると、以下のようになる（3次以上の微小項は無視する）。
\begin{equation}
\begin{aligned}
F(A+B)&=(A+B)+(A+B)^2 \\
&=A+B+(A^2+2AB+B^2) \\
&\approx A+B+A^2
\end{aligned}
\label{eq:appendix_welfare_correction_consistency_logic_example}
\end{equation}
ここで重要なのは、\textbf{「入力に含まれる2次の項\(B\)」と「関数によって生成される2次の項\(A^2\)」の両方が、計算結果として残る}という点である。

しかし、もしシミュレーションが1次近似で行われていたらどうなるだろうか。シミュレーション結果\(X^{sim}\)は、微細な\(B\)を無視して\(X^{sim}\approx A\)と出力される。これをそのまま非線形関数に代入してしまうと、
\begin{equation}
F(X^{sim})=A+A^2
\label{eq:appendix_welfare_correction_consistency_logic_inconsistent}
\end{equation}
となる。一見もっともらしい値が出るが、ここには重大な欠陥がある。本来足されるべきだった\textbf{「入力由来の2次の項\(B\)」が欠落している一方で、「関数由来の2次の項\(A^2\)」だけが計算されている}のである。同じ重要度を持つはずの要素のうち、片方だけを計算し、片方を無視するのは、計算として不整合であり、結果の信頼性を損なう。



本稿の分析においても、これと全く同じ問題が発生する。

\paragraph{①指数関数によるレベル変数の復元を行ってはならない理由}\label{sec:appendix_welfare_correction_consistency_p1}
シミュレーション結果として得られる対数乖離\(\hat{x}_t\)から、レベル変数\(x_t\)を復元する際、定義式である指数関数\(x_t=x_{ss}\exp(\hat{x}_t)\)を用いて計算してはいけない。
なぜなら、\textbf{非線形な指数関数をマクローリン展開すると}\(x_{ss}(1+\hat{x}_t+\frac{1}{2}\hat{x}_t^2+\dots)\)となり、\textbf{2次以上の項が含まれる}からである。入力である\(\hat{x}_t\)自体が2次の情報（上例の\(B\)）を欠落させているにも関わらず、計算式だけで2次以上の項を生成するのは、上述の不整合を引き起こす。

\paragraph{②線形復元した変数を非線形な効用関数に代入してはならない理由}\label{sec:appendix_welfare_correction_consistency_p2}
では、2次以上の項が出ないように\(x_t=x_{ss}(1+\hat{x}_t)\)と線形で復元すればよいかというと、それも不十分である。
もし、そうして復元した\(x_t\)（すなわち\(c_t,l_t\)）を、以下のような対数や二乗を含む非線形な効用関数
\begin{equation}
U(c_t^{H\to W},l_t^H)=\ln c_t^{H\to W}-\frac{\phi^H}{2}(l_t^H)^2
\label{eq:appendix_welfare_correction_consistency_utility_nonlinear}
\end{equation}
に直接代入してしまえば、結局は効用関数側が2次以上の項（\(A^2\)など）を生成することになり、入力情報の欠落（\(B\)の無視）との不整合が生じるからである。

\paragraph{③正しい対処法：効用関数自体の線形近似}\label{sec:appendix_welfare_correction_consistency_p3}
整合性を保つための正しい手順は、変数を代入する前に、まず効用関数\(U\)自体を定常状態近傍で1次近似（線形化）し、その式に出てくる\(\hat{c}_t,\hat{l}_t\)に、シミュレーションで得られた1次近似値\(\hat{c}_t,\hat{l}_t\)を代入することである。
このように関数を線形化しておけば、「線形」対「線形」の対応となり、整合性が保たれる。

\paragraph{④価格分散コスト\(\Delta\)の手動補正}\label{sec:appendix_welfare_correction_consistency_p4}
ただし、この線形化の手順をとると、本来は2次以上の項として評価されるべき重要な要素、すなわち「価格分散コスト\(\Delta_t\)」が式から消滅してしまう（1次近似の世界では\(\Delta_t\approx0\)となるため）。
そこで本稿では、この消えてしまったコスト\(\Delta_t\)を、\textcite{Woodford2003}にならい別途手計算で導出し、線形近似された厚生の値から事後的に差し引くという補正を行う。これにより、シミュレーションの整合性を保つとともに、価格のばらつきによる経済的損失を正確に評価に反映させることが可能となる。

次節より、その具体的な導出過程を示す。

\section{価格分散コストの補正式の導出}
\label{sec:appendix_welfare_correction_derivation}

補正の根拠は、線形近似モデルにおいてゼロとなる価格分散の影響を、効用関数の計算に復元することにある。以下にその導出過程を示す。

\paragraph{1.効用関数の線形近似}\label{sec:appendix_welfare_correction_derivation_p1}

まず、前節で示した期間効用関数\(U(c_t^{H\to W},l_t^H)\)を再掲する。
\begin{equation}
    U(c_t^{H\to W},l_t^H)=\ln c_t^{H\to W}-\frac{\phi^H}{2}(l_t^H)^2
\label{eq:appendix_welfare_correction_derivation_utility_restated}
\end{equation}
この関数の値が、定常状態\((c_{ss}^{H\to W},l_{ss}^H)\)からどの程度乖離するかを、1次テイラー展開（線形近似）を用いて評価する。
多変数関数の近似公式を適用すると、効用の変動部分は以下のように記述できる。
\begin{align}
    U(c_t^{H\to W},l_t^H)-U(c_{ss}^{H\to W},l_{ss}^H)&\approx\frac{\partial U}{\partial c_t^{H\to W}}(c_{ss}^{H\to W},l_{ss}^H)\cdot(c_t^{H\to W}-c_{ss}^{H\to W}) \label{eq:appendix_welfare_correction_derivation_utility_taylor_step1} \\
    &\quad+\frac{\partial U}{\partial l_t^H}(c_{ss}^{H\to W},l_{ss}^H)\cdot(l_t^H-l_{ss}^H)
\label{eq:appendix_welfare_correction_derivation_utility_taylor}
\end{align}
ここで、定常状態における各偏微分係数は以下の通りである。
\begin{align}
    \frac{\partial U}{\partial c_t^{H\to W}}(c_{ss}^{H\to W},l_{ss}^H)&=\frac{1}{c_{ss}^{H\to W}} \label{eq:appendix_welfare_correction_derivation_diff_c} \\
    \frac{\partial U}{\partial l_t^H}(c_{ss}^{H\to W},l_{ss}^H)&=-\phi^Hl_{ss}^H \label{eq:appendix_welfare_correction_derivation_diff_l}
\end{align}
また、ここでの変数\(\hat{x}_t\)を定常状態からの乖離\(\hat{x}_t=\frac{x_t-x_{ss}}{x_{ss}}\)と定義すると、レベル変数の乖離は\(x_t-x_{ss}=x_{ss}\hat{x}_t\)と変換できる。
これらを上式に代入して整理すると、以下の線形近似された効用の変動式が得られる。
\begin{align}
    U(c_t^{H\to W},l_t^H)-U(c_{ss}^{H\to W},l_{ss}^H)&\approx\frac{1}{c_{ss}^{H\to W}}\cdot(c_{ss}^{H\to W}\hat{c}_t^{H\to W})+(-\phi^Hl_{ss}^H)\cdot(l_{ss}^H\hat{l}_t^H) \label{eq:appendix_welfare_correction_derivation_utility_base_step} \\
    &=\hat{c}_t^{H\to W}-\phi^H(l_{ss}^H)^2\hat{l}_t^H
\label{eq:appendix_welfare_correction_derivation_utility_base}
\end{align}

\paragraph{2.労働投入量の近似と補正項の導出}\label{sec:appendix_welfare_correction_derivation_p2}

次に、労働投入量\(\hat{l}_t^H\)を生産関数から導出する。
第3章で示した集計生産関数\(y_t^H=a_t^Hl_t^H/\Delta_t^H\)を、まず労働投入量\(l_t^H\)について解く。
\begin{equation}
    l_t^H=\frac{y_t^H\Delta_t^H}{a_t^H}
\label{eq:appendix_welfare_correction_derivation_labor_def}
\end{equation}
この式の両辺について定常状態からの乖離（ハット変数）をとると、以下の関係が得られる。
\begin{equation}
    \hat{l}_t^H=\hat{y}_t^H-\hat{a}_t^H+\hat{\Delta}_t^H
\label{eq:appendix_welfare_correction_derivation_labor_approx}
\end{equation}

この式\eqref{eq:appendix_welfare_correction_derivation_labor_approx}を、そのまま式\eqref{eq:appendix_welfare_correction_derivation_utility_base}に代入すると、以下の式が得られる。
\begin{align}
    U(c_t^{H\to W},l_t^H)-U(c_{ss}^{H\to W},l_{ss}^H)&\approx\hat{c}_t^{H\to W}-\phi^H(l_{ss}^H)^2(\hat{y}_t^H-\hat{a}_t^H+\hat{\Delta}_t^H) \label{eq:appendix_welfare_correction_derivation_utility_combined_step1} \\
    &=\hat{c}_t^{H\to W}-\phi^H(l_{ss}^H)^2(\hat{y}_t^H-\hat{a}_t^H)-\phi^H(l_{ss}^H)^2\hat{\Delta}_t^H
\label{eq:appendix_welfare_correction_derivation_utility_combined}
\end{align}

シミュレーションにおいては、この第3項\(-\phi^H(l_{ss}^H)^2\hat{\Delta}_t^H\)が重要な役割を果たす。標準的な線形近似モデルではこの項が欠落してしまうため、本分析では特別に2次近似を用いて\(\hat{\Delta}_t^H\)を導出し、手動で厚生計算に反映させる。

以下に、その導出過程を記述する。

\paragraph{価格分散の関数の定義と近似方針}\label{sec:appendix_welfare_correction_derivation_disp_p1}

まず、第3章で導出した\textbf{物価指数の動学}（式\ref{eq:final_price_index_dynamics_H}）と\textbf{価格分散の動学}（式\ref{eq:final_dispersion_dynamics_H}）を出発点とする。ここで議論を簡潔にするため人口を\(N=1\)とする。
\begin{equation}
    1=(1-\xi^H)\left(\frac{\widetilde{p}_t^H}{p_t^H}\right)^{1-\theta^H}+\xi^H(\pi_t^H)^{\theta^H-1}
\label{eq:appendix_welfare_correction_derivation_pi_dynamics}
\end{equation}
\begin{equation}
    \Delta_t^H=(1-\xi^H)\left(\frac{\widetilde{p}_t^H}{p_t^H}\right)^{-\theta^H}+\xi^H(\pi_t^H)^{\theta^H}\Delta_{t-1}^H
\label{eq:appendix_welfare_correction_derivation_delta_dynamics}
\end{equation}
式\eqref{eq:appendix_welfare_correction_derivation_pi_dynamics}を相対価格\(\widetilde{p}_t^H/p_t^H\)について解き、それを式\eqref{eq:appendix_welfare_correction_derivation_delta_dynamics}に代入して相対価格を消去すると、価格分散\(\Delta_t^H\)はインフレ率\(\pi_t^H\)と前期の分散\(\Delta_{t-1}^H\)のみの関数として、以下のように定義できる。
\begin{equation}
    \Delta_t^H=(1-\xi^H)^{\frac{1}{1-\theta^H}}\left(1-\xi^H(\pi_t^H)^{\theta^H-1}\right)^{\frac{\theta^H}{\theta^H-1}}+\xi^H(\pi_t^H)^{\theta^H}\Delta_{t-1}^H
\label{eq:appendix_welfare_correction_derivation_delta_combined}
\end{equation}

この関数を、定常状態\((\pi_{ss}^H,\Delta_{ss}^H)=(1,1)\)の周りで2次テイラー展開（近似）する。
一般に、2変数関数\(f(x,y)\)の点\((x_0,y_0)\)周りでの2次近似公式は以下のように記述される。
\begin{align}
    f(x,y)&\approx f(x_0,y_0)+\frac{\partial f}{\partial x}(x_0,y_0)(x-x_0)+\frac{\partial f}{\partial y}(x_0,y_0)(y-y_0) \label{eq:appendix_welfare_correction_derivation_2nd_taylor_step1} \\
    &\quad+\frac{1}{2}\left[\frac{\partial^2 f}{\partial x^2}(x_0,y_0)(x-x_0)^2+2\frac{\partial^2 f}{\partial x\partial y}(x_0,y_0)(x-x_0)(y-y_0)+\frac{\partial^2 f}{\partial y^2}(x_0,y_0)(y-y_0)^2\right]
\label{eq:appendix_welfare_correction_derivation_2nd_taylor}
\end{align}
この公式を本モデルに適用する。ここで、インフレ率と価格分散の定常状態値は\(\pi_{ss}^H=1,\Delta_{ss}^H=1\)であるため、レベル変数の乖離\(x_t-1\)は、対数乖離\(\hat{x}_t\)と以下のように一致する。
\begin{equation}
    x_t-1=\frac{x_t-1}{1}=\frac{x_t-x_{ss}}{x_{ss}}=\hat{x}_t
\label{eq:appendix_welfare_correction_derivation_hat_identity}
\end{equation}
また、変数\(\pi_t^H\)と\(\Delta_{t-1}^H\)は第2項において分離した形（積の形）で入っており、第2項は\(\Delta_{t-1}^H\)について1次式であるため、\(\Delta_{t-1}^H\)に関する2階微分はゼロとなる。また交差項\(\hat{\pi}_t^H\hat{\Delta}_{t-1}^H\)は3次の微小量となるため無視できる。

したがって、求める近似式は以下の形式となる。
\begin{equation}
    \hat{\Delta}_t^H\approx\frac{\partial\Delta_t^H}{\partial\Delta_{t-1}^H}(1,1)\hat{\Delta}_{t-1}^H+\frac{\partial\Delta_t^H}{\partial\pi_t^H}(1,1)\hat{\pi}_t^H+\frac{1}{2}\frac{\partial^2\Delta_t^H}{\partial(\pi_t^H)^2}(1,1)(\hat{\pi}_t^H)^2
\label{eq:appendix_welfare_correction_derivation_delta_structure}
\end{equation}

\paragraph{偏微分係数の導出}\label{sec:appendix_welfare_correction_derivation_diff_p1}

まず、式\eqref{eq:appendix_welfare_correction_derivation_delta_combined}に基づき、全ての偏導関数を計算する。

\textbf{1.\(\Delta_{t-1}^H\)に関する1階偏微分} \\
第2項のみを微分する。
\begin{equation}
    \frac{\partial\Delta_t^H}{\partial\Delta_{t-1}^H}=\xi^H(\pi_t^H)^{\theta^H}
\label{eq:appendix_welfare_correction_derivation_partial_delta_lag}
\end{equation}

\textbf{2.\(\pi_t^H\)に関する1階偏微分} \\
合成関数の微分公式を用いて計算する。
\begin{align}
    \frac{\partial\Delta_t^H}{\partial\pi_t^H}&=(1-\xi^H)^{\frac{1}{1-\theta^H}}\frac{\theta^H}{\theta^H-1}\left(1-\xi^H(\pi_t^H)^{\theta^H-1}\right)^{\frac{\theta^H}{\theta^H-1}-1}\left(-\xi^H(\theta^H-1)(\pi_t^H)^{\theta^H-2}\right) \label{eq:appendix_welfare_correction_derivation_pi_first_diff_step1} \\
    &\quad+\xi^H\theta^H(\pi_t^H)^{\theta^H-1}\Delta_{t-1}^H \label{eq:appendix_welfare_correction_derivation_pi_first_diff_step2} \\
    &=-\theta^H\xi^H(1-\xi^H)^{\frac{1}{1-\theta^H}}\left(1-\xi^H(\pi_t^H)^{\theta^H-1}\right)^{\frac{1}{\theta^H-1}}(\pi_t^H)^{\theta^H-2} \label{eq:appendix_welfare_correction_derivation_pi_first_diff_step3} \\
    &\quad+\xi^H\theta^H(\pi_t^H)^{\theta^H-1}\Delta_{t-1}^H
\label{eq:appendix_welfare_correction_derivation_pi_first_diff}
\end{align}

\textbf{3.\(\pi_t^H\)に関する2階偏微分} \\
上記の1階偏微分をさらにもう一度\(\pi_t^H\)で微分する。第1項には積の微分公式を適用する。
\begin{align}
    \frac{\partial^2\Delta_t^H}{\partial(\pi_t^H)^2}&=-\theta^H\xi^H(1-\xi^H)^{\frac{1}{1-\theta^H}}\Bigg[\frac{1}{\theta^H-1}\left(1-\xi^H(\pi_t^H)^{\theta^H-1}\right)^{\frac{1}{\theta^H-1}-1}(-\xi^H(\theta^H-1)(\pi_t^H)^{\theta^H-2})\cdot(\pi_t^H)^{\theta^H-2} \label{eq:appendix_welfare_correction_derivation_pi_second_diff_step1} \\
    &\quad+\left(1-\xi^H(\pi_t^H)^{\theta^H-1}\right)^{\frac{1}{\theta^H-1}}\cdot(\theta^H-2)(\pi_t^H)^{\theta^H-3}\Bigg] \label{eq:appendix_welfare_correction_derivation_pi_second_diff_step2} \\
    &\quad+\xi^H\theta^H(\theta^H-1)(\pi_t^H)^{\theta^H-2}\Delta_{t-1}^H
\label{eq:appendix_welfare_correction_derivation_pi_second_diff}
\end{align}

次に、計算したこれらの偏導関数を定常状態\((\pi_{ss}^H,\Delta_{ss}^H)=(1,1)\)で評価する。

\textbf{1.ラグ項の係数の評価}
\begin{equation}
    \frac{\partial\Delta_t^H}{\partial\Delta_{t-1}^H}(1,1)=\xi^H(1)^{\theta^H}=\xi^H
\label{eq:appendix_welfare_correction_derivation_eval_lag}
\end{equation}

\textbf{2.1次の係数の評価}
\begin{align}
    \frac{\partial\Delta_t^H}{\partial\pi_t^H}(1,1)&=-\theta^H\xi^H(1-\xi^H)^{\frac{1}{1-\theta^H}}\left(1-\xi^H\right)^{\frac{1}{\theta^H-1}}(1)+\xi^H\theta^H(1)(1) \label{eq:appendix_welfare_correction_derivation_eval_pi_first_step1} \\
    &=-\theta^H\xi^H(1-\xi^H)^{\frac{1}{1-\theta^H}+\frac{1}{\theta^H-1}}+\xi^H\theta^H \label{eq:appendix_welfare_correction_derivation_eval_pi_first_step2} \\
    &=-\theta^H\xi^H(1-\xi^H)^0+\xi^H\theta^H \label{eq:appendix_welfare_correction_derivation_eval_pi_first_step3} \\
    &=-\theta^H\xi^H+\theta^H\xi^H=0
\label{eq:appendix_welfare_correction_derivation_eval_pi_first}
\end{align}

\textbf{3.2次の係数の評価}
\begin{align}
    \frac{\partial^2\Delta_t^H}{\partial(\pi_t^H)^2}(1,1)&=-\theta^H\xi^H(1-\xi^H)^{\frac{1}{1-\theta^H}}\Bigg[-\xi^H(1-\xi^H)^{\frac{1}{\theta^H-1}-1}(1)+(1-\xi^H)^{\frac{1}{\theta^H-1}}(\theta^H-2)\Bigg] \label{eq:appendix_welfare_correction_derivation_eval_pi_second_step1} \\
    &\quad+\xi^H\theta^H(\theta^H-1)(1)(1) \label{eq:appendix_welfare_correction_derivation_eval_pi_second_step2} \\
    &=\theta^H(\xi^H)^2(1-\xi^H)^{-1}-\theta^H\xi^H(\theta^H-2)(1)+\xi^H\theta^H(\theta^H-1) \label{eq:appendix_welfare_correction_derivation_eval_pi_second_step3} \\
    &=\frac{\theta^H(\xi^H)^2}{1-\xi^H}+\theta^H\xi^H\left[-(\theta^H-2)+(\theta^H-1)\right] \label{eq:appendix_welfare_correction_derivation_eval_pi_second_step4} \\
    &=\frac{\theta^H(\xi^H)^2}{1-\xi^H}+\theta^H\xi^H \label{eq:appendix_welfare_correction_derivation_eval_pi_second_step5} \\
    &=\theta^H\xi^H\left(\frac{\xi^H}{1-\xi^H}+1\right)=\frac{\theta^H\xi^H}{1-\xi^H}
\label{eq:appendix_welfare_correction_derivation_eval_pi_second}
\end{align}

\paragraph{結論：価格分散の2次近似式と線形モデルでの含意}\label{sec:appendix_welfare_correction_derivation_conclusion_p1}

以上の計算結果を近似式に代入することで、以下の最終的な関係式が得られる。
\begin{align}
    \hat{\Delta}_t^H&\approx\xi^H\hat{\Delta}_{t-1}^H+0\cdot\hat{\pi}_t^H+\frac{1}{2}\left(\frac{\theta^H\xi^H}{1-\xi^H}\right)(\hat{\pi}_t^H)^2 \label{eq:appendix_welfare_correction_derivation_delta_approx_step} \\
    &=\xi^H\hat{\Delta}_{t-1}^H+\frac{\theta^H\xi^H}{2(1-\xi^H)}(\hat{\pi}_t^H)^2
\label{eq:appendix_welfare_correction_derivation_delta_approx}
\end{align}
これが最終的な価格分散の2次近似式である。

この結果（式\ref{eq:appendix_welfare_correction_derivation_delta_approx}）は、線形近似モデルにおいて価格分散項がゼロとなることの数学的な証明となっている。具体的には、上式の導出過程で確認した通り、インフレ率\(\hat{\pi}_t^H\)に関する1次の係数（偏微分係数）は計算の結果ちょうど0となる。したがって、もし方程式系全体を1次のオーダーで近似（線形化）した場合、2次以上の微小量である右辺第2項（\((\hat{\pi}_t^H)^2\)を含む項）は無視され、式は以下のように退化する。
\begin{equation}
    \hat{\Delta}_t^H\approx\xi^H\hat{\Delta}_{t-1}^H
\label{eq:appendix_welfare_correction_derivation_delta_degenerated}
\end{equation}
定常状態から出発する場合、初期値は\(\hat{\Delta}_0^H=0\)であるため、この再帰式に従えば、将来にわたり常に\(\hat{\Delta}_t^H=0\)となる。これが、標準的な線形近似シミュレーションにおいて価格分散による厚生ロスが消失してしまう理由である。

\paragraph{厚生計算に用いる最終的な効用関数}\label{sec:appendix_welfare_correction_derivation_final_p1}

最後に、導出された価格分散の近似式\eqref{eq:appendix_welfare_correction_derivation_delta_approx}を、先述の効用の近似式\eqref{eq:appendix_welfare_correction_derivation_utility_combined}に代入することで、本分析のプログラムにおいて実際に計算される最終的な効用関数の式が得られる。

\begin{equation}
    U(c_t^{H\to W},l_t^H)-U(c_{ss}^{H\to W},l_{ss}^H)\approx\hat{c}_t^{H\to W}-\phi^H(l_{ss}^H)^2(\hat{y}_t^H-\hat{a}_t^H)-\phi^H(l_{ss}^H)^2\left(\xi^H\hat{\Delta}_{t-1}^H+\frac{\theta^H\xi^H}{2(1-\xi^H)}(\hat{\pi}_t^H)^2\right)
\label{eq:appendix_welfare_correction_derivation_final_utility}
\end{equation}     % 厚生補正（ 価格分散の動学 ）
% !TeX root = ../../main.tex
% sections/app/appendix_insulation_effect.tex

\chapter{遮断効果を生みだす方程式系の導出}
\label{chap:appendix_insulation_effect}

本付録では、自国でショックが起こった際に自国金融政策が外国に関連する以下の変数に影響を与えないことを証明する。
\(\{\lambda_t^{H/*},\lambda_t^{F/*},p_t^{F*},i_t^F,c_t^{H\to F},c_t^{F\to F},y_t^F,l_t^F,\bar{p}_t^{F*},\pi_t^{F*},\widetilde{p}_t^{F*},v_t^F,w_t^F,t_t^{F*},b_t^H,\Delta_t^F\}\)
このような現象は\textcite{ColeObstfeld1991}の遮断効果として知られている。
遮断効果の要因は上記の変数が外国の方程式群のみによって決定されることによる。
そこでまずこれらの外国の方程式を列挙し、それらのみによって上記の変数が決定されることを示す。

議論の見通しを良くするため証明で用いる変数を定義しておく。
\begin{itemize}
    \item\label{itm:appendix_insulation_effect_def_lambda_star} \textbf{自国の外国通貨所得の限界効用：}\(\lambda_t^{H/*}\equiv e_t\lambda_t^H\)\\
    これは自国の外国通貨所得の限界効用、すなわち1単位の外国通貨で消費することにより自国代表的個人が得られる効用を意味する。
\end{itemize}

\section{外国のコア方程式系と独立性の証明}
\label{sec:appendix_insulation_effect_core}

まず、外国の主要な内生変数を決定するコア方程式系を導出する。



\paragraph{1.外国財市場の均衡式}
\label{sec:appendix_insulation_effect_core_market_clearing}

まず、外国財市場の実質需給均衡式を出発点とする。
\begin{equation}
y_t^F=c_t^{F\to F}+\frac{N}{M}c_t^{H\to F}
\label{eq:appendix_insulation_effect_core_market_clearing_real}
\end{equation}
この両辺に外国財価格\(p_t^{F*}\)を乗じ、外国通貨建ての名目額に変形する。
\begin{equation}
    p_t^{F*}y_t^F=p_t^{F*}c_t^{F\to F}+\frac{N}{M}p_t^{F*}c_t^{H\to F}
    \label{eq:appendix_insulation_effect_core_market_clearing_nom_step}
\end{equation}
右辺に代表的家計の反映需要関数を代入する。さらに、消費の最適化条件より成立する名目総消費支出の関係式\(p_t^{F\to W*}c_t^{F\to W}=1/\lambda_t^{F/*}\)および\(p_t^{H\to W}c_t^{H\to W}=1/\lambda_t^H\)を適用し、自国の外国通貨所得の限界効用\(\lambda_t^{H/*}\equiv e_t\lambda_t^H\)を用いて一気に整理すると、以下のようになる。

\begin{align}
p_t^{F*}y_t^F&=p_t^{F*}\left(\alpha^F\frac{p_t^{F\to W*}}{p_t^{F*}}c_t^{F\to W}\right)+\frac{N}{M}p_t^{F*}\left((1-\alpha^H)\frac{p_t^{H\to W}}{e_tp_t^{F*}}c_t^{H\to W}\right) \label{eq:appendix_insulation_effect_core_market_clearing_nom_demand_step1} \\
&=\alpha^F(p_t^{F\to W*}c_t^{F\to W})+\frac{N}{M}\frac{1-\alpha^H}{e_t}(p_t^{H\to W}c_t^{H\to W}) \label{eq:appendix_insulation_effect_core_market_clearing_nom_demand_step2} \\
&=\alpha^F\frac{1}{\lambda_t^{F/*}}+\frac{N}{M}\frac{1-\alpha^H}{e_t\lambda_t^H} \label{eq:appendix_insulation_effect_core_market_clearing_nom_demand_step3} \\
&=\alpha^F\frac{1}{\lambda_t^{F/*}}+\frac{N}{M}\frac{1-\alpha^H}{\lambda_t^{H/*}}
\label{eq:appendix_insulation_effect_core_market_clearing_nom_demand_final}
\end{align}

これにより、以下の外国名目総生産の決定式が得られる。
\begin{equation}
    p_t^{F*}y_t^F=\alpha^F\frac{1}{\lambda_t^{F/*}}+\frac{N}{M}\frac{1-\alpha^H}{\lambda_t^{H/*}}
    \label{eq:appendix_insulation_effect_core_market_clearing_nom_final}
\end{equation}

\paragraph{2.外国のオイラー方程式}
\label{sec:appendix_insulation_effect_core_euler}

外国家計の異時点間の最適化条件より、
\begin{equation}
\lambda_t^{F/*}=\beta_t^F(1+i_t^F)E_t[\lambda_{t+1}^{F/*}]
\label{eq:appendix_insulation_effect_core_euler_final}
\end{equation}

\paragraph{3.外国のフィリップス曲線群}
\label{sec:appendix_insulation_effect_core_phillips}

外国企業の価格設定行動を記述する非線形方程式群である。

まず、外国最適価格\(\widetilde{p}_t^{F*}\)の決定式：
\begin{equation}
    (\widetilde{p}_t^{F*})^{1+\theta^F}=\frac{\theta^F}{\theta^F-1}\frac{v_t^F}{w_t^F}
    \label{eq:appendix_insulation_effect_core_phillips_opt_p}
\end{equation}
次に、補助変数\(v_t^F\)の再帰式：
\begin{equation}
    v_t^F=\phi^F\frac{(y_t^F)^2\Delta_t^F}{(a_t^F)^2}(p_t^{F*})^{2\theta^F}+\beta_t^F\xi^FE_t[v_{t+1}^F]
    \label{eq:appendix_insulation_effect_core_phillips_v_recursive}
\end{equation}
続いて、補助変数\(w_t^F\)の再帰式：
\begin{equation}
    w_t^F=\lambda_t^{F/*}(1-\tau_t^F)y_t^F(p_t^{F*})^{\theta^F}+\beta_t^F\xi^FE_t[w_{t+1}^F]
    \label{eq:appendix_insulation_effect_core_phillips_w_recursive}
\end{equation}
そして価格分散\(\Delta_t^F\)の動学式：
\begin{equation}
    \Delta_t^F=(1-\xi^F)M\left(\frac{\widetilde{p}_t^{F*}}{p_t^{F*}}\right)^{-\theta^F}+\xi^F\left(\frac{p_t^{F*}}{p_{t-1}^{F*}}\right)^{\theta^F}\Delta_{t-1}^F
    \label{eq:appendix_insulation_effect_core_phillips_delta_dynamics}
\end{equation}
そして物価指数の動学：
\begin{equation}
    (p_t^{F*})^{1-\theta^F}=(1-\xi^F)M(\widetilde{p}_t^{F*})^{1-\theta^F}+\xi^F(p_{t-1}^{F*})^{1-\theta^F}
    \label{eq:appendix_insulation_effect_core_phillips_price_dynamics}
\end{equation}

\paragraph{4.外国の金融政策式}
\label{sec:appendix_insulation_effect_core_policy}

外国中央銀行の政策ルールにおいて、\(\pi_t^{F*}\)と\(\pi_{ss}^F\)を\(p_t^{F*}\)と\(p_{ss}^{F*}\)で展開すると以下が得られる。
\begin{equation}
    i_t^F=i_{ss}^F+\phi_{\pi}^F\left(\frac{p_t^{F*}}{p_{t-1}^{F*}}-1\right)+\varepsilon_t^{i,F}
    \label{eq:appendix_insulation_effect_core_policy_rule}
\end{equation}

\paragraph{5.国際債券に関する修正オイラー方程式}
\label{sec:appendix_insulation_effect_core_modified_uip}

市場均衡式\eqref{eq:appendix_insulation_effect_core_market_clearing_nom_final}に現れた自国の外国通貨所得の限界効用\(\lambda_t^{H/*}\)がどのように決定されるかを示す。
自国代表的家計の国際債券に関するオイラー方程式を出発点とする。
\begin{equation}
    \lambda_t^He_t=(1+i_t^F)\beta_t^HE_t[\lambda_{t+1}^He_{t+1}]
    \tag{\ref{eq:uip_condition_home}}
\end{equation}
この左辺は定義により\(\lambda_t^{H/*}\)であり、右辺の期待値の中身は\(\lambda_{t+1}^{H/*}\)である。
したがって以下の式が得られる。
\begin{equation}
    \lambda_t^{H/*}=\beta_t^H(1+i_t^F)E_t[\lambda_{t+1}^{H/*}]
    \label{eq:appendix_insulation_effect_uip_final}
\end{equation}

この方程式系に含まれる変数のうち\(\lambda_t^{F/*},v_t^F,w_t^F,\lambda_t^{H/*}\)の4つは
将来の期待値に依存して現在の値が決定されるジャンプ変数（前向きな変数）である。
具体的には外国のオイラー方程式\eqref{eq:appendix_insulation_effect_core_euler_final}
のように次期の期待値が含まれることで現在の行動が規定される性質をもつ。
対照的に、価格分散\(\Delta_t^F\)のように過去の履歴に拘束される変数は状態変数と呼ばれる。

9つの外国関連変数はこれら9本の方程式系において決定される。
そこであらためてこの9本の方程式系に現れる変数をすべて列挙すると、まず9つの外国関連変数
\(\{\lambda_t^{H/*},\lambda_t^{F/*},p_t^{F*},i_t^F,\widetilde{p}_t^{F*},v_t^F,w_t^F,y_t^F,\Delta_t^F\}\)
パラメータ
\(\{\alpha^H,\alpha^F,N,M, \theta^F,\xi^F,\phi^F,i_{ss}^F,\phi_{\pi}^F,\pi_{ss}^F\}\)
そして外生ショック
\(\{\beta_t^H,\beta_t^F,a_t^F, \tau_t^F,\varepsilon_t^{i,F}\}\)
である。
したがって9つの外国関連変数は\(\beta_t^H\)などの影響は受けるが
自国の金融政策\(i_t^H\)や自国の生産者物価指数\(p_t^H\)、自国の生産\(y_t^H\)といった
自国の内生変数の影響は受けないことがわかる。
このように9本の方程式系が閉じていることが遮断効果を生んでいる。

\section{その他外国関連変数の決定}
\label{sec:appendix_insulation_effect_other}
上記で決定された9つの外国関連変数
\(\{\lambda_t^{H/*},\lambda_t^{F/*},p_t^{F*},i_t^F,\widetilde{p}_t^{F*},v_t^F,w_t^F,y_t^F,\Delta_t^F\}\)
よりさらに残りの外国関連変数が決定される。

\paragraph{6.自国の外国財消費\(c_t^{H\to F}\)の決定}
\label{sec:appendix_insulation_effect_other_chf}

次に自国家計による外国財の消費\(c_t^{H\to F}\)（外国にとっては輸出需要）を考える。
自国家計の消費の最適化条件より、自国通貨建ての名目総消費支出は、自国の所得の限界効用の逆数と等しくなる。
\begin{equation}
p_t^{H\to W}c_t^{H\to W}=\frac{1}{\lambda_t^H}
\label{eq:appendix_insulation_effect_other_chf_budget}
\end{equation}
一方、代表的家計の需要関数（式\eqref{eq:model_demand_HF}）の両辺に\(p_t^F\)を乗じると、自国の外国財への名目支出額（自国通貨建て）が得られる。
\begin{equation}
p_t^Fc_t^{H\to F}=(1-\alpha^H)(p_t^{H\to W}c_t^{H\to W})=\frac{1-\alpha^H}{\lambda_t^H}
\label{eq:appendix_insulation_effect_other_chf_nom_exp}
\end{equation}
一物一価の法則\(p_t^F=e_tp_t^{F*}\)を用いて左辺を書き換え、さらに両辺を\(e_t\)で割ることで外国通貨建てに変換する。
\begin{equation}
p_t^{F*}c_t^{H\to F}=\frac{1-\alpha^H}{e_t\lambda_t^H}
\label{eq:appendix_insulation_effect_other_chf_currency_step}
\end{equation}
ここで自国の外国通貨所得の限界効用の定義\(\lambda_t^{H/*}\equiv e_t\lambda_t^H\)を用いると、右辺は連結変数のみで表現できる。
\begin{equation}
p_t^{F*}c_t^{H\to F}=\frac{1-\alpha^H}{\lambda_t^{H/*}}
\label{eq:appendix_insulation_effect_other_chf_currency_final}
\end{equation}
この式を\(c_t^{H\to F}\)について解くと、以下の決定式が得られる。
\begin{equation}
    c_t^{H\to F}=\frac{1-\alpha^H}{p_t^{F*}\lambda_t^{H/*}}
    \label{eq:appendix_insulation_effect_other_chf_final}
\end{equation}
右辺の\(p_t^{F*}\)はコア方程式系から、\(\lambda_t^{H/*}\)は自国の外国通貨所得の限界効用の決定式\eqref{eq:appendix_insulation_effect_uip_final}から、それぞれ自国政策とは独立に決定済みである。したがって、自国の外国財消費\(c_t^{H\to F}\)もまた、自国金融政策の影響を受けない。

\paragraph{7.外国の外国財消費\(c_t^{F\to F}\)の決定}
\label{sec:appendix_insulation_effect_other_cff}

まず、外国家計による外国財の消費\(c_t^{F\to F}\)を考える。
外国家計の消費の最適化条件より、外国通貨建ての名目総消費支出は、所得の限界効用の逆数と等しくなる。
\begin{equation}
p_t^{F\to W*}c_t^{F\to W}=\frac{1}{\lambda_t^{F/*}}
\label{eq:appendix_insulation_effect_other_cff_budget}
\end{equation}
一方、代表的家計の需要関数（式\eqref{eq:model_demand_FF}）の両辺に\(p_t^{F*}\)を乗じると、外国財への名目支出額が得られる。
\begin{equation}
p_t^{F*}c_t^{F\to F}=\alpha^F(p_t^{F\to W*}c_t^{F\to W})
\label{eq:appendix_insulation_effect_other_cff_nom_exp}
\end{equation}
これら2式を結合すると、外国財への名目支出額は、所得の限界効用のみによって決定されることがわかる。
\begin{equation}
p_t^{F*}c_t^{F\to F}=\frac{\alpha^F}{\lambda_t^{F/*}}
\label{eq:appendix_insulation_effect_other_cff_combined}
\end{equation}
この式を\(c_t^{F\to F}\)について解くと、以下の決定式が得られる。
\begin{equation}
    c_t^{F\to F}=\frac{\alpha^F}{p_t^{F*}\lambda_t^{F/*}}
    \label{eq:appendix_insulation_effect_other_cff_final}
\end{equation}
右辺の\(p_t^{F*}\)および\(\lambda_t^{F/*}\)は、コアとなる方程式系において既に自国政策から独立して決定された変数である。したがって、外国財消費\(c_t^{F\to F}\)もまた、自国金融政策の影響を受けずに一意に定まる。

\paragraph{8.外国の労働\(l_t^F\)の決定}
\label{sec:appendix_insulation_effect_other_labor}

続いて、外国の労働を考える。
個別生産関数の集計式を\(l_t^F\)について逆算することで、以下の決定式が得られる。
\begin{equation}
    l_t^F=\frac{y_t^F\Delta_t^F}{a_t^F}
    \label{eq:appendix_insulation_effect_other_labor_final}
\end{equation}
ここで、\(y_t^F\)（生産）と\(\Delta_t^F\)（価格分散）はコア方程式系で決定済みの変数であり、\(a_t^F\)は外生変数である。したがって、労働\(l_t^F\)も独立して決定される。

\paragraph{9.外国の正規化物価指数の決定}
\label{sec:appendix_insulation_effect_other_ppi}

外国の正規化された物価指数(PPI)\(\bar{p}_t^{F*}\)は、正規化された物価指数と通常の物価指数の関係式\(p_t^{F*}=M^{\frac{1}{1-\theta^F}}\bar{p}_t^{F*}\)より決定される。
\begin{equation}
    \bar{p}_t^{F*}=M^{-\frac{1}{1-\theta^F}}p_t^{F*}
    \label{eq:appendix_insulation_effect_other_ppi_final}
\end{equation}
なお、右辺の\(p_t^{F*}\)はコア方程式系ですでに決定されている。

\paragraph{10.外国のインフレ率の決定}
\label{sec:appendix_insulation_effect_other_inflation}

外国のグロス・インフレ率\(\pi_t^{F*}\)は、定義より以下のように決定される。
\begin{equation}
\pi_t^{F*}=\frac{\bar{p}_t^{F*}}{\bar{p}_{t-1}^{F*}}
\label{eq:appendix_insulation_effect_other_inflation_def}
\end{equation}

\paragraph{11.外国の一括移転\(t_t^{F*}\)の決定}
\label{sec:appendix_insulation_effect_other_transfer}

最後に、外国政府から家計への一括移転\(t_t^{F*}\)を考える。
外国政府の予算制約式より、移転額は名目GDPと税率によって決定される。
\begin{equation}
    t_t^{F*}=\tau_t^Fp_t^{F*}y_t^F
    \label{eq:appendix_insulation_effect_other_transfer_final}
\end{equation}
\(p_t^{F*}\)と\(y_t^F\)は決定済みであり、税率\(\tau_t^F\)は外生的な確率過程に従う。したがって、一括移転\(t_t^{F*}\)も自国政策から独立している。

\paragraph{12.自国対外純資産\(b_t^H\)の決定と独立性}
\label{sec:appendix_insulation_effect_other_nfa}

最後に、自国対外純資産\(b_t^H\)の決定について考察する。
自国の資源制約式は以下の通りである。
\begin{equation}
    p_t^{H\to W}c_t^{H\to W}+b_{t+1}^H=p_t^Hy_t^H+(1+i_{t-1}^F)\frac{e_t}{e_{t-1}}b_t^H
    \label{eq:appendix_insulation_effect_other_nfa_budget_raw}
\end{equation}
この式の両辺に、自国の所得の限界効用\(\lambda_t^H(=1/(p_t^{H\to W}c_t^{H\to W}))\)を乗じる。左辺第1項は1となり、右辺の資産収益項は定義\(\lambda_t^{H/*}\equiv e_t\lambda_t^H\)を用いて変形できる。
\begin{align}
1+\lambda_t^Hb_{t+1}^H&=\lambda_t^Hp_t^Hy_t^H+(1+i_{t-1}^F)\frac{e_t\lambda_t^H}{e_{t-1}}b_t^H \label{eq:appendix_insulation_effect_other_nfa_transform_step1} \\
&=\lambda_t^Hp_t^Hy_t^H+(1+i_{t-1}^F)\frac{\lambda_t^{H/*}}{e_{t-1}}b_t^H
\label{eq:appendix_insulation_effect_other_nfa_transform_step2}
\end{align}
ここで、右辺第2項を、前期の\(\lambda_{t-1}^{H/*}\equiv e_{t-1}\lambda_{t-1}^H\)を用いて書き換えると、
\begin{equation}
\frac{\lambda_t^{H/*}}{e_{t-1}}b_t^H=\frac{\lambda_t^{H/*}}{\lambda_{t-1}^{H/*}}(\lambda_{t-1}^Hb_t^H)
\label{eq:appendix_insulation_effect_other_nfa_shift}
\end{equation}
となる。

また、右辺第1項の自国名目所得（効用評価額）\(\lambda_t^Hp_t^Hy_t^H\)については、コブ＝ダグラス型選好の下では以下の手順で導出される。
\begin{align}
\lambda_t^Hp_t^Hy_t^H&=\lambda_t^Hp_t^H(c_t^{H\to H}+\frac{M}{N}c_t^{F\to H}) \label{eq:appendix_insulation_effect_other_nfa_income_step1} \\
&=\lambda_t^Hp_t^H\left(\alpha^H\frac{p_t^{H\to W}}{p_t^H}c_t^{H\to W}\right)+\lambda_t^Hp_t^H\frac{M}{N}\left((1-\alpha^F)\frac{p_t^{F\to W*}}{p_t^{H*}}c_t^{F\to W}\right) \label{eq:appendix_insulation_effect_other_nfa_income_step2} \\
&=\lambda_t^H\alpha^H(p_t^{H\to W}c_t^{H\to W})+\lambda_t^Hp_t^H\frac{M}{N}(1-\alpha^F)\frac{p_t^{F\to W*}}{p_t^H/e_t}c_t^{F\to W} \label{eq:appendix_insulation_effect_other_nfa_income_step3} \\
&=\alpha^H+\frac{M}{N}(1-\alpha^F)e_t\lambda_t^H(p_t^{F\to W*}c_t^{F\to W}) \label{eq:appendix_insulation_effect_other_nfa_income_step4} \\
&=\alpha^H+\frac{M}{N}(1-\alpha^F)\lambda_t^{H/*}\frac{1}{\lambda_t^{F/*}} \label{eq:appendix_insulation_effect_other_nfa_income_step5} \\
&=\alpha^H+\frac{M}{N}(1-\alpha^F)\frac{\lambda_t^{H/*}}{\lambda_t^{F/*}}
\label{eq:appendix_insulation_effect_other_nfa_income_final}
\end{align}

これらを合わせると、貯蓄の効用価値を表す変数\(X_t\equiv\lambda_t^Hb_{t+1}^H\)に関する以下の差分方程式が得られる。
\begin{equation}
    1+X_t=\left[\alpha^H+\frac{M}{N}(1-\alpha^F)\frac{\lambda_t^{H/*}}{\lambda_t^{F/*}}\right]+\left[(1+i_{t-1}^F)\frac{\lambda_t^{H/*}}{\lambda_{t-1}^{H/*}}\right]X_{t-1}
    \label{eq:appendix_insulation_effect_other_nfa_diff_eq}
\end{equation}
この式\eqref{eq:appendix_insulation_effect_other_nfa_diff_eq}の右辺にあるすべての項は、前項までに証明した通り、すべて自国金融政策\(i^H\)から独立して決定される変数である。したがって、\(X_t\)の経路は自国金融政策の影響を受けない。

さらに、定常状態近傍での線形近似を考えると、定常状態で\(b_{ss}^H=0\)であるため、\(X_{ss}=0\)となる。
変数\(X_t=\lambda_t^Hb_{t+1}^H\)を定常状態周りで1次近似（テイラー展開）すると、
\begin{equation}
X_t\approx X_{ss}+\lambda_{ss}^H(b_{t+1}^H-b_{ss}^H)+b_{ss}^H(\lambda_t^H-\lambda_{ss}^H)
\label{eq:appendix_insulation_effect_other_nfa_approx_step}
\end{equation}
ここで\(X_{ss}=0,b_{ss}^H=0\)を代入すると、
\begin{equation}
X_t\approx\lambda_{ss}^Hb_{t+1}^H
\label{eq:appendix_insulation_effect_other_nfa_linear_approx}
\end{equation}
という関係が得られる。\(\lambda_{ss}^H\)は正の定数であるため、\(X_t\)が政策不変であれば、自国対外純資産\(b_{t+1}^H\)の経路もまた、自国金融政策の影響を受けないことが証明される。

このように、外国経済のすべての実質変数および名目変数は、自国の金融政策変数\(i_t^H\)を含む方程式系から構造的に切り離されており、影響を受けないことが確認できる。      % 遮断効果の数学的証明